% manuscript.tex — AFM 合体稿骨格 (Phase W, 2026-08-03)
% Title/abstract: 裁定済み 2026-08-03 (SKELETON.md §1–§2, md5 6e8758a6 時点)
% Source of truth: repo rd-quartics-classification, frozen tag v1.0.4 (commit 6c123af).
% The Lean artifact wins on any discrepancy. Mathematics is FROZEN; only
% arrangement and narration move (SKELETON §5).
% Line-pin convention: blob/v1.0.4/<path>#L<n>.
% Forbidden phrases (grep must be 0 before any freeze; DISCIPLINE §2):
%   mutatis mutandis / we believe / trivially / machine-checked / grevlex /
%   analogous combinatorial expansion / standard binomial identity / obvious /
%   clearly / easy to see / automatic
% (use "kernel-verified" / "kernel-certified", never "machine-checked")

\documentclass[11pt]{amsart}
\usepackage{amsmath,amssymb,amsthm}
\usepackage[hidelinks]{hyperref}
\hypersetup{pdftitle={Rational-derived quartics in Lean: a corrected
repeated-root classification and the boundary of Conjecture 1},
pdfauthor={Hiroki Fukui},
pdfsubject={Rational-derived polynomials, elliptic curves, formal
verification in Lean 4}}
\usepackage{microtype}
\usepackage{array}
\usepackage{booktabs}
\usepackage{listings}
% Engine split for Lean Unicode in listings: under XeTeX (tectonic) the
% literate mechanism does not fire for non-ASCII, so use a system
% monospace font with the required glyphs; under pdfLaTeX keep literate.
\ifdefined\XeTeXrevision
  \usepackage[no-math]{fontspec}
  \newfontfamily\leanmono{Menlo}[Scale=MatchLowercase]
  \lstset{basicstyle=\leanmono\footnotesize}
\else
  \lstset{basicstyle=\ttfamily\footnotesize}
\fi
\lstset{columns=fullflexible, keepspaces=true, escapeinside={(*@}{@*)},
  literate={ℚ}{{\ensuremath{\mathbb{Q}}}}1 {∀}{{\ensuremath{\forall}}}1
    {∃}{{\ensuremath{\exists}}}1 {∧}{{\ensuremath{\wedge}}}1
    {≠}{{\ensuremath{\neq}}}1 {→}{{\ensuremath{\to}}}1
    {⟨}{{\ensuremath{\langle}}}1 {⟩}{{\ensuremath{\rangle}}}1
    {σ}{{\ensuremath{\sigma}}}1 {γ}{{\ensuremath{\gamma}}}1
    {ρ}{{\ensuremath{\rho}}}1 {ζ}{{\ensuremath{\zeta}}}1
    {'}{{\textquotesingle}}1}

\newtheorem{theorem}{Theorem}[section]
\newtheorem{lemma}[theorem]{Lemma}
\newtheorem{proposition}[theorem]{Proposition}
\newtheorem{remark}[theorem]{Remark}
\newtheorem{definition}[theorem]{Definition}

\newcommand{\Q}{\mathbb{Q}}
\newcommand{\Z}{\mathbb{Z}}
\newcommand{\R}{\mathbb{R}}

% ---- artifact reference macros (all pins to the frozen tag) ----
\newcommand{\repourl}{https://github.com/hirokifukui/rd-quartics-classification}
\newcommand{\repotag}{v1.1.0}
\newcommand{\repocommit}{a5cf446}
\newcommand{\conceptdoi}{10.5281/zenodo.21766032}
\newcommand{\blueprinturl}{https://hirokifukui.github.io/rd-quartics-classification/blueprint/}
% \leanfile{path} -> hyperlinked path at the frozen tag
\newcommand{\leanfile}[1]{\href{\repourl/blob/\repotag/#1}{\texttt{#1}}}
% \leanline{path}{line}{name} -> declaration name pinned to file+line
\newcommand{\leanline}[3]{\href{\repourl/blob/\repotag/#1\#L#2}{\texttt{#3}}}

\title[Rational-derived quartics in Lean]{Rational-derived quartics in
Lean:\\ a corrected repeated-root classification\\ and the boundary of
Conjecture~1}

\author{Hiroki Fukui}
\address{Research Institute of Criminal Psychiatry / Sex Offender Medical
Center, Tokyo, Japan; and Department of Neuropsychiatry, Kyoto University,
Kyoto, Japan. ORCID: 0009-0008-7122-522X}
\email{fukui@somec.org}

% MSC: 候補 68V20 + 11D41 + 11G05 — 凍結前に三重確認 (SKELETON §4;
% Manthe の front-matter 事故に学ぶ)
\subjclass[2020]{Primary 68V20; Secondary 11G05, 11D41}
\keywords{rational-derived polynomials, elliptic curves, formal
verification, Lean, mathlib, proof auditing, proof repair}

\begin{document}

\begin{abstract}
% 裁定済み v2 (2026-08-03). 結語一文 = release 結語と同文・不動.
A polynomial over $\Q$ is \emph{rational-derived} if it and all of its
derivatives split completely over $\Q$.  Buchholz and MacDougall
(J.~Number Theory 81 (2000)) presented a classification of
rational-derived quartics, their Theorem~7, whose four-distinct-root
clause is entangled with their Conjecture~1.  We report a three-part
study of this classification, conducted through a Lean~4 formalization.
First, an \emph{audit}: the published proof of Theorem~7 contains a
circularity and a quantifier gap --- the argument assumes that the
critical vertical translate is rational-derived, an assumption
equivalent to the splitting it is required to prove (the equivalence is
kernel-certified), and then concludes universally without discharging
it; a single quartic, read over two base fields, refutes over $\Q$
the implication left once the circular assumption is removed, and over
$\Q(\sqrt{3441})$ the statement's natural analogue under the full
$K$-derived hypothesis.  Second, a \emph{repair}: Theorem~$7'$, a
rank-free image theorem identifying the normalised $p(2,1,1)$ sector
exactly with the image of the rational points of the elliptic curve
$E\colon z^{2}=w(w-6)(w+18)$ (Cremona \texttt{576i2}) under an
explicit rational map --- which, together with a kernel-verified affine
normalisation and the elementary multiplicity cases, yields an
unconditional classification of all repeated-root quartics; both
directions reduce to explicit polynomial
certificates together with elementary field and order arguments, and
\emph{no rank computation enters}
--- a proof found because the library lacked Mordell--Weil rank.  The
main statement compiles with axiom footprint exactly
\texttt{[propext, Classical.choice, Quot.sound]}.  Third, a
\emph{boundary}: a kernel-verified bridge identifies the normalised
form of Conjecture~1 with the nonexistence of nondegenerate rational
points on an explicit splitting surface; the three double covers
governing the vertical-translate route are proved nontrivial over the
splitting extension, with the kernel boundary of each claim marked in
the text, and the fibre correspondence between the disjuncts of
Theorem~7 and those covers is closed in the kernel.  Conjecture~1
itself remains open: this development maps its boundary, it does not
resolve it.  The formalization detected the gap, sharpened the
statements, and certified the repair; we record the mathematics
together with the lessons of the formalization.
\end{abstract}

\maketitle

%% ================================================================
%% §1 INTRODUCTION (~4–5p; 最後に書く — SKELETON §8 執筆順序 6)
%% SOURCE: 新規 + jnt Intro + RNT Intro + READING_NOTES 総合章 §D
%% ================================================================
\section{Introduction}\label{sec:intro}

Around 1950, M.~S.~Klamkin exhibited a quartic polynomial,
$(x-193)(x-141)(x+167)^2$, with the property that it and all of its
derivatives have integer roots.  For decades it stood as a curiosity:
Carroll, writing in the Monthly in 1989, conjectured that it was
essentially the only quartic of its kind with three or more distinct
roots \cite{Carroll1989}; in the same year Richard Guy's problem column
printed four pairs of roots of further examples, each pair larger than
the last, the fourth running to twenty digits \cite{Guy1989}.  In 1995
Buchholz and Kelly organized the family by an elliptic curve
\cite{BK1995}, and in 2000 Buchholz and MacDougall presented a
classification of all such \emph{rational-derived} quartics, their
Theorem~7, reducing what remained to a conjecture \cite{BM2000}.  One
outcome of the present study (Theorem~\ref{thm:curve}) is that the
scattered record has a single organizing point: Guy's four printed
pairs are, in order and up to one root exchange, the images of the
consecutive multiples $2P$, $3P$, $4P$, $5P$ of a single point $P$ on a
single elliptic curve of conductor $576$ --- with Klamkin's quartic the
first of them.  The examples, scattered over a seventy-six-year record, can be
identified --- across the different models used in the literature ---
with consecutive multiples of a single point on a single curve.

Call $f\in\Q[x]$ \emph{rational-derived} if $f$ and every derivative of
$f$ split completely over $\Q$; for a field $K\supseteq\Q$, call $f$
\emph{$K$-derived} if $f$ and all its positive-degree derivatives split
completely over $K$, and \emph{properly} $K$-derived if it is
$K$-derived but not rational-derived.  Linear, quadratic and cubic
rational-derived polynomials are completely understood; for quartics,
sorting by root multiplicities, the types $p(4)$ and $p(3,1)$ are
always rational-derived, the type $p(2,2)$ never is, the type
$p(2,1,1)$ --- one double root, two simple --- is the substantial
positive theory, and the type $p(1,1,1,1)$ is the subject of
Conjecture~1 of \cite{BM2000}, which asserts that no quartic with four
distinct roots is rational-derived and which remains open.  (For
quintics with a triple root the corresponding nonexistence was settled
by Flynn \cite{Flynn2001}.)  Theorem~7 of \cite{BM2000} presents the
full classification, and subsequent constructions
\cite{BC2016,Carrillo2019} cite it as the classification of the
quartic case.

\subsection{What this paper does}\label{sec:intro:parts}

This paper reports a three-part study of that classification, conducted
through a Lean~4 \cite{Lean} formalization, and its results sort the quartic
picture into three strata.  \emph{The repeated-root stratum is now
classified unconditionally}: Part~II states and proves Theorem~$7'$ in
image form --- the normalised $p(2,1,1)$ sector is exactly the image of
the rational points of the curve \texttt{576i2} under an explicit map,
with no subgroup or coset restriction --- closes the statement-level
links (gate form against natural form, normalised family against
arbitrary split quartics) in the kernel, and assembles the
unconditional classification of all repeated-root quartics
(Theorem~\ref{thm:assembled}).  \emph{The four-distinct-root stratum is mapped to its
boundary}: Part~I locates, by a unit-by-unit transcription of the
published proof, the exact step of \cite{BM2000} that is not
established --- a printed assumption equivalent to the conclusion it
should deliver, followed by a universal claim the preceding analysis
does not cover --- and refutes the reconstructed implication with a
single quartic read over two base fields; Part~III then proves, in the
kernel, that Conjecture~1 is equivalent to the nonexistence of
nondegenerate rational points on an explicit splitting surface, and
shows that all three routes the printed argument could take are
governed by double covers that are nontrivial over the splitting
extension.  \emph{Conjecture~1 itself remains open}: nothing in this
paper resolves it, and \S\ref{sec:open} displays it, deliberately
unproved, in the exact form it takes in the development.  Along the
way the audit corrects two Mordell--Weil ranks in the sources
(\S\ref{sec:ranks}) and yields the historical unification above.

Throughout, ``the formalization'' names the whole program --- the
verbatim transcription and audit of the published text, the
computer-algebra gating of every statement-level definition, the statement-level
design pressure, and the kernel verification --- not the kernel alone.
The arc of the paper is that this program \emph{detects, sharpens, and
certifies}: the audit detected a gap that twenty-six years of use had
not surfaced; the discipline of writing formal statements sharpened
loose ones (the coset restriction of the 1995 model, the excluded
parameters $a\in\{0,1\}$, the degenerate loci of the $a$-map, all now
exactly delimited); and the kernel certified the repair,
unconditionally, with the standard three-axiom footprint
(\S\ref{sec:formal:axioms}).  Attribution is kept
at that resolution deliberately: it would be inaccurate to say simply
that ``Lean found the gap.''

\subsection{What we learned}\label{sec:lessons}

The lesson we consider most transferable is that a limitation of the
library produced a better theorem.  The published route to Theorem~7
runs through the Mordell--Weil rank of the curve; mathlib \cite{mathlib},
at the pinned commit, did not provide the Mordell--Weil rank machinery
required by the published route, so the conditional
first version of our development carried the rank as a disclosed axiom.
The pressure to remove it forced a different question --- what do the
two gate witnesses alone determine? --- and the answer
(\S\ref{sec:repair:forward}) is an explicit rational reconstruction of
a point of the curve from the witnesses, discovered by Gr\"obner
elimination and certified by the kernel, which proves the completeness
direction with \emph{no rank input at all}.  The repaired theorem is
not merely formalized; it is unconditional in a way Theorem~7's own statement was not (the
coset-form classification of \cite{BK1995}, in the repeated-root
sector, was itself unconditional; the conditionality sat in
Theorem~7's exhaustion claim), and the disclosed axiom, deliberately retained,
marks in the kernel's own per-declaration report the point where the
condition was removed (\S\ref{sec:formal:axioms}).

This experience composes three lessons already recorded in the
Annals of Formalized Mathematics.  Loeffler and Stoll describe how a formal statement of the
Riemann hypothesis nearly slipped at its edge --- a junk value that,
had it been zero, would have made the displayed statement a false
proposition superficially resembling RH --- and warn that
statement-level slips are the failure mode to fear
\cite{LoefflerStoll}; the present paper is the same failure type
located in the informal literature, where it appears to have gone
unnoted through twenty-six years of citation, and the statement-fidelity gates of
\S\ref{sec:fidelity} are our systematic response.  Best, Birkbeck,
Brasca and their coauthors describe formalizability driving the choice
of proof, back to older and more elementary arguments
\cite{BrascaEtAl}; here the drive went further, to a rank-free proof
absent from the cited record.  Ludwig and Merten describe formalization changing
the mathematics --- generalizations forced by the formal setting
\cite{LudwigMerten}; here the change is the unconditional two-gate
form of the statement itself.
The type of the present paper --- the audit, repair and certification
of a published, peer-reviewed classification --- combines the three in
a single audit-and-repair arc: the slip at the
statement's edge is its subject, the return to the sound 1995 kernel
its method, and the sharpened, rank-free theorem its outcome.

\subsection{The artifact}\label{sec:intro:artifact}

Every kernel-verified statement cited in this paper is hyperlinked, at
theorem level, to a specific declaration at the frozen release
\repotag{} (commit \texttt{\repocommit}) of \url{\repourl};
load-bearing computer-algebra claims are indexed to their scripts and
captured logs (\S\ref{sec:formal:artifact}); the dependency graph
with per-node kernel status is the blueprint at \url{\blueprinturl},
and the archived snapshot carries concept DOI \texttt{\conceptdoi}.
The live development comprises 4{,}347 lines of Lean across twelve
files in two components --- \texttt{Gap/} (the audit and the boundary)
and \texttt{Thm7Prime/} (the repair) --- with 209 theorems, one
disclosed axiom off the main line, and no \texttt{sorry}; toolchain,
pins, continuous verification and quantitative details are
\S\ref{sec:formal}.  Part~I is \S\S\ref{sec:gap}--\ref{sec:ranks};
Part~II is \S\S\ref{sec:repair}--\ref{sec:fidelity}; Part~III is
\S\ref{sec:boundary}; \S\ref{sec:formal} records the formalization and
the division of labour, and \S\ref{sec:open} the open sector.

%% ================================================================
%% PART I — THE AUDIT
%% ================================================================
\part*{Part I. The audit}

%% ----------------------------------------------------------------
%% §2 THE GAP (~5p)
%% SOURCE: jnt §2 ほぼ全量 + blueprint Ch1–2 (v1.0.4)
%% 四段・非感情的: reconstructed statement → formal consequence →
%% counterexample → replacement
%% ----------------------------------------------------------------
\section{The gap}\label{sec:gap}

The object of this part is a single section of \cite{BM2000}:
\S2.2.4, ``Vertical translation of a quartic,'' which supplies the last
inferential step of that paper's Theorem~7.  The audit is a survey of
one section, not a verdict on the paper: the identities established
there are exact and are re-proved in the kernel below, and
typographical errors and inferential gaps are distinct phenomena, kept
distinct throughout.  For the audit the section was segmented into
eleven inferential units $\mathrm{T0}$--$\mathrm{T10}$ covering it with
no residue; each unit is transcribed as a Lean object in
\leanfile{Gap/BMThm7Transcript.lean}, over a general field (linearly
ordered where an order is needed), so that a published $K$-derived
quartic can act as a transparent negative control below.  Line
references $\ell.n$ are to the archived working layout of the text
(hash in \texttt{Gap/PROVENANCE.md}); that layout derives from the
authors' 2009 revision, whose \S2.2.4 differs from the published
version only typographically, and all quotations below are given in the
published wording.

\subsection{The promise and the assumption, quoted}\label{sec:gap:quotes}

The section opens with its promise ($\ell$.529--534):

\begin{quote}\small
``If we now allow our rational-derived quartics to undergo a vertical
translate by a rational distance, so that the (highest) local minimum,
$x_{\min}$ say, is moved up to become a double root, then the remaining
two roots, $r$ and $s$ say, of the quartic could possibly lie in a
quadratic extension of $\Q$. Certainly, we have no a priori reason to
expect the extra roots to be rational (see Fig. 2). None-the-less, in
this section we show that the latter is precisely the case.''
\end{quote}

\noindent This is the terminus of the transcript: the proposition the
section undertakes to prove, transcribed as
\leanline{Gap/BMThm7Transcript.lean}{201}{T0\_claim}.  The argument
then names its standing datum.  Having introduced
$f(x)=x(x-1)(x-a)(x-b)$ and a rational translate amount $c$, it writes
($\ell$.553--558):

\begin{quote}\small
``We assume that the resulting quartic, $F(x)$ say, given by
\[
F(x) = x(x-1)(x-a)(x-b) + c \tag{4}
\]
is rational-derived and has a double root. As before, we require
$\Delta(F)=0$\dots''
\end{quote}

\noindent (The equation is displayed there with tag (4); nothing else
is elided.)  The verb is \emph{assume}; the assumption is transcribed
as \leanline{Gap/BMThm7Transcript.lean}{83}{T2\_assumption}.  And,
having excluded the degenerate cases and established the discriminant
identities recalled below, the section concludes verbatim
($\ell$.615--620):

\begin{quote}\small
``Previously one may have thought that the class of rational derived
quartics with four distinct roots split into two types: those
obtainable by a rational vertical translate from a p(2,1,1) quartic and
those not so obtainable. The result of all the previous work shows that
the latter class is empty and so any rational-derived quartic with four
distinct roots can be vertically translated into a rational-derived
quartic with a double root.''
\end{quote}

\subsection{The formal consequence: a circle and a quantifier jump}
\label{sec:gap:circle}

The content of the printed assumption is fixed by a one-line
observation.

\begin{lemma}[\leanline{Gap/BMThm7Transcript.lean}{89}{KDerived\_translate\_iff}]
\label{lem:circ}
Let $x_0\in K$ be a critical point of the quartic $f\in K[x]$, and put
$g(x)=f(x)-f(x_0)$.  Vertical translation leaves every positive-degree
derivative unchanged.  Consequently, if $f$ is $K$-derived, then $g$ is
$K$-derived if and only if $g$ itself splits over $K$ (the pinned
kernel statement; equivalently, since $(x-x_0)^2$ always splits, if and
only if the residual quadratic $q$ in $g(x)=(x-x_0)^2 q(x)$ splits over
$K$).
\end{lemma}

Under the standing hypothesis of the section, the printed assumption is
therefore \emph{equivalent} to the missing splitting assertion for the
selected critical translate --- and in particular it already implies
the existential conclusion asserted at the end of the section.  The
pointwise equivalence is kernel-certified
(\leanline{Gap/BMThm7Transcript.lean}{106}{T2\_iff\_conclusion}).  The
structure of the argument is thus: (1) translate at the selected
critical point (the highest local minimum); (2) assume the translate is
rational-derived --- by Lemma~\ref{lem:circ}, assume the required
splitting for that translate; (3) analyse, under that assumption, the
resulting repeated-root family; (4) conclude, in the sentence quoted
above, that the class not covered by the assumption ``is empty'',
without an argument.  The gap is a circularity together with a
quantifier jump: the exhaustion is asserted for all rational-derived
quartics, while the preceding analysis covers exactly those satisfying
the assumption.  The kernel transcript records this structurally: no
lemma of the transcribed chain takes the splitting of $f''$ as a
hypothesis --- the full derived condition enters only through the
assumption, which is precisely the pointwise splitting assertion
required for the selected translate --- and the concluding universal
statement is derivable from the chain precisely when that assumption is
supplied (\leanline{Gap/BMThm7Transcript.lean}{316}{T9\_via\_T2}).  The
implication the section actually asserts --- from the hypotheses
consumed by the transcribed chain to the splitting of the translate ---
is transcribed, \emph{as a definition and not a theorem}, as
\leanline{Gap/BMThm7Transcript.lean}{308}{T9\_bridge}; it is refuted in
\S\ref{sec:gap:onequartic} below.

The origin of the slip is instructive.  For \emph{cubics} the analogous
bridging step is a theorem: after forcing a double root, a single root
remains, and its rationality is free from the rationality of the root
sum --- this is exactly the argument of \S2.1 of \cite{BM2000}, which
is sound.  For quartics two roots remain, the root sum fixes only their
sum, and the published argument carries the cubic conclusion across by
analogy.

The correspondence between the published passages and the transcribed
chain is summarized in the following table; the axiom footprint of
every node is the standard three (\S\ref{sec:formal}).

\begin{center}
\footnotesize
\setlength{\tabcolsep}{4pt}
\begin{tabular}{@{}>{\raggedright\arraybackslash}p{0.16\textwidth}>{\raggedright\arraybackslash}p{0.24\textwidth}>{\raggedright\arraybackslash}p{0.19\textwidth}>{\raggedright\arraybackslash}p{0.20\textwidth}>{\raggedright\arraybackslash}p{0.11\textwidth}@{}}
\toprule
Published passage & Formal node & Hypotheses consumed & Output & Role \\
\midrule
assumption preceding (4) &
\leanline{Gap/BMThm7Transcript.lean}{106}{T2\_iff\_conclusion} &
quartic is $K$-derived; selected critical point $x_0$ &
printed assumption $\Leftrightarrow$ splitting of the residual quadratic at $x_0$ &
missing bridge \\[2pt]
discriminant identities &
\leanline{Gap/BMThm7Transcript.lean}{157}{deltaF\_eq\_sym\_mul\_deltaFp\_cubed};
\leanline{Gap/BMThm7Transcript.lean}{258}{D6\_eq\_sq\_of\_vieta},
\leanline{Gap/BMThm7Transcript.lean}{278}{D6\_nonneg\_of\_vieta} &
root and critical-point identities (Vieta data) &
the displayed identities; $D_6\ge 0$ &
non-circular chain \\[2pt]
concluding paragraph &
\leanline{Gap/BMThm7Transcript.lean}{316}{T9\_via\_T2} &
the printed assumption &
splitting at the selected translate &
quantifier jump \\
\bottomrule
\end{tabular}
\end{center}

\subsection{One quartic over two base fields}\label{sec:gap:onequartic}

\begin{theorem}[One quartic, two base fields]\label{thm:gap}
Let
\[
f(x)=x(x-136)(x-52)(x+312),\qquad
F(X)=448^{-4}\,f(448X-312),
\]
so that
$F(X)=X(X-1)\bigl(X-\tfrac{39}{56}\bigr)\bigl(X-\tfrac{13}{16}\bigr)$,
and let $K=\Q(\sqrt{3441})$.  Then:

\textup{(a)} over $\Q$, $f$ has four distinct rational roots, $f'$
splits, $f$ avoids every symmetric configuration excluded in
{\cite[\S2.2.4]{BM2000}}, and no critical vertical translate of $f$
splits over $\Q$; hence the algebraic hypotheses actually used by the
non-circular portion of the argument do not imply its concluding claim
(\leanline{Gap/BMThm7Transcript.lean}{489}{not\_T9\_bridge\_Q});

\textup{(b)} $F$ is properly $K$-derived with four distinct roots, and
no critical vertical translate of $F$ splits over $K$; hence the
natural $K$-analogue of the bridging claim is false
(\leanline{Gap/BMThm7Transcript.lean}{977}{not\_T0\_claim\_K} at the
critical point selected by the published argument);

\textup{(c)} $F$ is the affine normalisation of the $K$-derived quartic
exhibited by Stroeker {\cite[Ex.~5.2]{Stroeker2006}} at $t=-7/17$.
\end{theorem}

The identity
$F(X)=448^{-4}\,f(448X-312)$
carries the roots $0,136,52,-312$ of $f$ to
$\tfrac{39}{56},1,\tfrac{13}{16},0$; the two parts of the theorem are
two readings of this single affine class.  (Different normalisations of
the class appear in the artifact: dividing the roots of $f$ by $136$
gives parameters $(13/34,\,-39/17)$ with selected critical point
$x_0=13/17$, the coordinates used in the kernel certificates; the
normalisation $(39/56,\,13/16)$ used in the text is the one matching
Stroeker's example.  Throughout, ``highest local minimum'' abbreviates
the transcript's algebraic second-derivative critical-point predicate,
stated over a linearly ordered field; over $\R$, under the
nondegeneracy hypotheses in force, it coincides with the higher of the
two local minima.)

\emph{Over $\Q$.}  The roots $0,136,52,-312$ of $f$ are distinct
rationals, and
\[
f'(x)=4x^3+372x^2-103168x+2206464=4(x+221)(x-24)(x-104)
\]
splits over $\Q$.  In the normalised coordinates $(a,b)=(39/56,\,13/16)$,
\[
a-b-1=-\tfrac{125}{112},\qquad a-b+1=\tfrac{99}{112},\qquad
a+b-1=\tfrac{57}{112},
\]
all nonzero, so $f$ avoids every symmetric configuration excluded in
the argument, and the discriminant identities of
\S\ref{sec:gap:circle} hold for it as identities.  Since $f'(x_0)=0$,
each $f(x)-f(x_0)$ has a double root at $x_0$; dividing by $(x-x_0)^2$
leaves a quadratic $q_{x_0}$:
\[
\begin{array}{llll}
x_0=-221: & q(x)=x^2-318x+40131, & \operatorname{disc}=-59400&=-66\cdot 30^2;\\[2pt]
x_0=24: & q(x)=x^2+172x-43904, & \operatorname{disc}=205200&=57\cdot 60^2;\\[2pt]
x_0=104: & q(x)=x^2+332x+6656, & \operatorname{disc}=83600&=209\cdot 20^2.
\end{array}
\]
Neither $57$ nor $209$ is a square, and the first discriminant is
negative; hence no critical translation of $f$ has rational remaining
roots.  Each identity above is a polynomial identity over $\Z$ and can
be checked directly; the full chain (the factorization of $f'$, the
three divisions, the three discriminants, and the non-squareness
assertions) is certified by the kernel under the standard three axioms
(\leanfile{Gap/BMThm7Gap.lean}), unconditionally over $\Q$.  We note
that $f$ itself is not rational-derived:
$\operatorname{disc}(f'')=40^2\cdot 3441$.  This is exactly what part
(a) requires --- the counterexample must live outside the circular
assumption while satisfying all algebraic premises actually invoked
after its removal --- and it is what part (b) turns to advantage.

\emph{Over $K=\Q(\sqrt{3441})$.}  Over $K$ the second derivative also
splits:
\[
F''(X)=\frac{5376X^2-6744X+1859}{448},\qquad
\text{with roots}\ \frac{843\pm 5\sqrt{3441}}{1344}\in K,
\]
so $F$ is $K$-derived, and properly so, since it is not
rational-derived.  Its critical points are given by
\[
F'(X)=\frac{(4X-3)(14X-13)(64X-13)}{896},\qquad
X_0\in\Bigl\{\tfrac{13}{64},\ \tfrac34,\ \tfrac{13}{14}\Bigr\},
\]
and the residual discriminants of the three critical translates are
\[
-66\Bigl(\tfrac{15}{224}\Bigr)^2,\qquad
57\Bigl(\tfrac{15}{112}\Bigr)^2,\qquad
209\Bigl(\tfrac{5}{112}\Bigr)^2
\]
--- the same square classes as over $\Q$, as they must be, this being
the same quartic.  Now $K$ is real, so the first is not a square in
$K$; and a rational number $r$ is a square in $K$ only if $r$ or
$3441r$ is a square in $\Q$, which fails for $57$ and $209$ (neither
$57\cdot 3441=196137$ nor $209\cdot 3441=719169$ is a square).  Hence
\emph{no} critical vertical translate of $F$ splits over $K$, although
$F$ satisfies the full $K$-derived hypothesis.  The statement at the
critical point selected by the argument ($X_0=13/14$, the highest local
minimum, with residual quadratic
$X^2-\tfrac{73}{112}X+\tfrac{13}{6272}$) is kernel-certified over an
explicit subfield model of $K$
(\leanline{Gap/BMThm7Transcript.lean}{977}{not\_T0\_claim\_K}); the
failure at the remaining two critical points is by the direct
computation above.  Both refutations, over $\Q$ and over $K$, funnel
through the same square class $209$.

Two statements should be kept separate.
Theorem~\ref{thm:gap} does not assert that the statement of
{\cite[Thm.~7]{BM2000}} over $\Q$ is false --- a counterexample over
$\Q$ would refute Conjecture~1 of the same paper, which is open
(\S\ref{sec:open}).  What fails is the passage from the reduction to
the stated exhaustion; whether that passage holds over $\Q$ under the
full rational-derived hypothesis is precisely the unresolved
classification question, entangled with Conjecture~1.  Part (b) shows
that in the quadratic-field setting to which the paper extends the
theory, the analogous passage fails as such.  As a statement about the
literature: we found no independent proof or published repair in the
sources and citation chains examined through July 2026; the exhaustion
over $\Q$ must accordingly be regarded as unproved.\footnote{A
revised version of the paper, dated June 21, 2009, is available from
the first author's website
(\url{https://www.ralphthetriangle.com/papers/rational-derived-polynomials};
accessed July 2026); its \S2.2.4 is unchanged.}

\subsection{Provenance, the surviving parts, and the replacements}
\label{sec:gap:provenance}

The affine class of the counterexample is not new.  In
{\cite[Ex.~5.2]{Stroeker2006}} Stroeker exhibits proper $K$-derived
quartics with parameter $U(t)=(2t^2+2t-1)/(5t+4)$; at $t=-7/17$ (listed
there with $D=3\cdot 31\cdot 37=3441$) one has $U=-13/17$ and roots
$1,\,-51/13,\,-17/7,\,-3$, which the affine map $X=13(1-u)/64$ carries
to $0,\,1,\,39/56,\,13/16$ --- the quartic $F$ of
Theorem~\ref{thm:gap}.  What is new here is the observation that this
known $K$-derived quartic, read over $\Q$, satisfies all algebraic
premises actually invoked by the non-circular portion of
{\cite[\S2.2.4]{BM2000}} after the printed assumption is removed, while
falsifying the concluding implication of that portion; and that over
$K$ it refutes the natural $K$-analogue of the bridge outright.  A
witness to the gap in the proof of the family's central classification
was thus already present in the family's own literature.

The audit is a survey, not a demolition, and it closes with an
accounting of what survives and what was replaced.  The translate setup
is sound.  The displayed discriminant identities are exact: the identity
$\Delta(f')=-16\,D_6(a,b)$ at $\ell$.606 and the $-2^{12}$ form of
$\Delta F$ at $\ell$.568 hold as printed (here $\Delta$ denotes the
resultant-normalised convention of that paper, to be distinguished from
the ordinary polynomial discriminant, which satisfies
$\operatorname{disc}(f')=+4\,D_6$ --- kernel-proved at the level of the
closed formulas, with the coefficient bridge; the resultant-convention
identifications themselves are certified in both computer-algebra
systems [C]).  Within the kernel the closed forms are reconciled:
\leanline{Gap/BMThm7Transcript.lean}{157}{deltaF\_eq\_sym\_mul\_deltaFp\_cubed}
proves that the closed form of $\Delta F$ equals
$\mathrm{sym}\cdot(\text{closed form of }\Delta(f'))^3$ with prefactor
exactly $1$ --- a tower nowhere printed with that prefactor.  Of the two
mutually inconsistent printed prefactors ($-2^{12}$ against $D_6^3$ at
$\ell$.568 and $2^8$ against $\Delta(f')^3$ at $\ell$.611, which with
$\Delta(f')=-16D_6$ contradict one another), the reconciliation
adjudicates the $2^8$ as a typographical error, immaterial to the
argument.  Three further local
repairs are recorded, each replacing a printed argument by a stronger,
kernel-checked one.  First, the printed stationary-point analysis of
the surface $z=D_6(a,b)$ inspects three stationary points and two
sample values and cannot by itself establish global non-negativity on
an unbounded domain; the conclusion $D_6\ge 0$ is nonetheless exact and
is re-proved in analysis-free form, valid over any linearly ordered
field: on the Vieta locus $D_6$ is $64$ times a square
(\leanline{Gap/BMThm7Transcript.lean}{258}{D6\_eq\_sq\_of\_vieta},
\leanline{Gap/BMThm7Transcript.lean}{278}{D6\_nonneg\_of\_vieta}).
Second, the disposal of the three symmetric configurations, attributed
in the paper to a closed-access source, is replaced by a self-contained
proof: in each symmetric case $K$-derivedness forces a rational point
on the conic $r^2=3(p^2+q^2)$, which has only the trivial solution by
$3$-adic descent
(\leanline{Gap/BMThm7Transcript.lean}{636}{symF\_zero\_not\_KDerived}),
eliminating the external citation from the load path.  Third, a
``Conversely'' assertion in the same section does not follow from the
preceding identities; it is not used in the sequel, and no repair is
required.

Finally, the audit was conducted under a pre-registration discipline:
the segmentation into units, the classification of each unit, and seven
numbered predictions about where the argument would and would not
formalise were fixed in a document whose hash was frozen before any
Lean build was attempted, and the adjudication record postdates it in
full.  All seven predictions were confirmed; two were confirmed in a
form stronger than predicted.  Predictions and confirmations were both
produced within this project, so the record is presented as procedural
transparency, not as independent verification.

%% ----------------------------------------------------------------
%% §3 RANKS AND E₀ (~2.5p)
%% SOURCE: jnt §3–4
%% ----------------------------------------------------------------
\section{Two corrected ranks and the curve $E_0$}\label{sec:ranks}

The audit also recomputed the numerical curve data of \cite{BM2000} and
\cite{Stroeker2006}.  Two entries required correction; a single curve
turned out to organize the family's whole record.  The computations of
this section are retained as scripts and logs in the archived
development (\texttt{scripts/rank/}, a SageMath~10.8 audit; the rank of
$E$ itself is additionally certified in Magma,
\texttt{scripts/magma/}).  The rank statements of
Theorem~\ref{thm:ranks} are unconditional ($2$-descent; no analytic
input is consumed by any statement of this paper).

\begin{theorem}[Corrected ranks]\label{thm:ranks}
Let $E\colon z^2=w(w-6)(w+18)$ be the auxiliary curve of \cite{BM2000}.
\begin{itemize}
\item[(a)] $\operatorname{rk} E(\Q(\sqrt{33}))=2$, not $1$ as recorded
in {\cite[Table~5]{BM2000}}.
\item[(b)] The elliptic curve underlying Example~4.1 of
\cite{Stroeker2006} (identical in the earlier report
\cite{Stroeker2002}) has Mordell--Weil rank $1$ over $\Q$, not $2$ as
stated there.
\end{itemize}
\end{theorem}

\emph{(a) Table 5, $m=33$.}  With $E_m$ the quadratic twist of $E$ by
$m$ (for $m=33$, minimal model $y^2=x^3-x^2-2097x+28305$), and the
twist identity
$\operatorname{rk}E(\Q(\sqrt m))=
\operatorname{rk}E(\Q)+\operatorname{rk}E_m(\Q)$ used throughout
{\cite[\S4]{BM2000}}: unconditional two-descent gives
$\operatorname{rk}E(\Q)=1$ and $\operatorname{rk}E_{33}(\Q)=1$, the
latter with explicit generator $(-35,240)$ on the minimal model,
saturation index $1$ (the analytic rank agrees, as an auxiliary check);
hence $\operatorname{rk}E(\Q(\sqrt{33}))=2$.  A direct computation over
$\Q(\sqrt{33})$ independently returns rank $2$, with the two
independent generators $(-12,36)$ and $(3-3s,\,36-12s)$, $s=\sqrt{33}$,
on $E$.  The $m=33$ entry is the sole discrepancy we assert.  The row-by-row
audit retained in the artifact is partial and not load-bearing: seven
of the nine entries of Table~4 are confirmed there by unconditional
two-descent and the remaining two agree with the analytic ranks
reported in the source; of Table~5, the $m=73$ row is confirmed
unconditionally and the $m=57$ row analytically only; the sweep, including
its aborted run, is preserved as such in \texttt{scripts/rank/}.

\emph{(b) Example 4.1 of \cite{Stroeker2006}.}  The curve of that
example is the quartic model
\[
Z^2 \;=\; 3\,(t^4+4t^3+4t^2+3),
\]
with rational points (e.g.\ $(t,Z)=(0,\pm 3)$), whose Jacobian is
$\Q$-isomorphic to $E_0$ below; by the rank of the curve we mean the
Mordell--Weil rank of that Jacobian.  Both versions state that this
curve has rank $2$.  Unconditional two-descent gives rank bounds
$(1,1)$, hence the Mordell--Weil rank is $1$; the point $(-8,36)$
generates $E_0(\Q)$ modulo its torsion (saturation index $1$, checked
unconditionally; log archived with the release accompanying this
paper).  The analytic rank is also
$1$, an independent consistency check.  The example's substantive
conclusion is unaffected: the correspondence asserted there is
independent of the rank computation, and rank $1$ still makes the
Mordell--Weil group infinite, so the infinitude of examples is
preserved.  (The result first circulated as Econometric Institute
Report EI 2002-30 (2002) and was then published in condensed form as
\cite{Stroeker2006}; Example~4.1, with the rank-$2$ statement, appears
identically in both versions.  The correction applies to both; we
address it primarily to the published version.)

\begin{theorem}[One curve through the family]\label{thm:curve}
Let $E_0\colon y^2=x^3-156x+560$, of conductor $576$.  Then:
\begin{itemize}
\item[(a)] $E_0$ is $\Q$-isomorphic to the auxiliary curve
$E\colon z^2=w(w-6)(w+18)$ of \cite{BM2000}, so that Tables~4 and~5 of
that paper are rank tables for the quadratic twists of $E_0$;
\item[(b)] $E_0$ is $\Q$-isomorphic to the Jacobian of the quartic
model in Example~4.1 of \cite{Stroeker2006}; in particular the
corrected rank of Theorem~\ref{thm:ranks}(b) is the rank of $E_0(\Q)$
itself, which is $1$;
\item[(c)] let $P=(-12,36)$ on $E\colon z^2=w(w-6)(w+18)$, a point of
infinite order generating $E(\Q)$ modulo its torsion $(\Z/2\Z)^2$, and
let $a(w,z)$ be the map \eqref{eq:amap}.  Then
\[
\begin{aligned}
a(2P)&=\frac{90}{77}, &
a(3P)&=\frac{497610}{167167},\\[2pt]
a(4P)&=-\frac{128665027260}{69283212011}, &
a(5P)&=\frac{4173952380480465660}{7010366636418797651},
\end{aligned}
\]
and these are, in order and up to the root-exchange symmetry
$a\leftrightarrow 1/a$, the normalised parameters of the quartics
$x^2(x-p)(x-q)$ of the \emph{four} pairs $(p,q)$ printed in Guy's
problem column \cite{Guy1989} --- the fourth pair after exchanging the
two simple roots.  The first pair $(308,360)$ yields a translate of
Klamkin's quartic $(x-193)(x-141)(x+167)^2$ (circa 1950; see
\cite{Carroll1989}).  (Guy separately credited M.~J.~Miller with a
refutation of Carroll's uniqueness conjecture; Miller's example is not
printed there, and we make no identification of it.)
\end{itemize}
\end{theorem}

For (a) no computation is needed: the substitution $w=x-4$ carries
$z^2=w(w-6)(w+18)$ to $y^2=(x-4)(x-10)(x+14)=x^3-156x+560$, an identity
the reader can check directly; both curves have conductor $576$ ($E_0$
is \texttt{576i2} in Cremona's tables --- the very curve of
Theorem~$7'$).  For (b), the quartic $Z^2=g(t)$ with
$g(t)=3t^4+12t^3+12t^2+9$ has classical invariants
$I=12ae-3bd+c^2=468$ and $J=72ace+9bcd-27ad^2-27b^2e-2c^3=-15120$
(writing $g=at^4+bt^3+ct^2+dt+e$; for this convention and the Jacobian
$y^2=x^3-27Ix-27J$ see, e.g., \cite{Cremona1997}), hence Jacobian
$y^2=x^3-12636x+408240$; setting $x=9X$, $y=27Y$ and dividing by $3^6$
gives $Y^2=X^3-156X+560$, i.e.\ $E_0$.  Both steps can be checked by
hand.  For (c): translating Klamkin's quartic
$p(x)=(x-193)(x-141)(x+167)^2$ --- whose derivative has the integer
roots $\{169,-2,-167\}$ and whose second derivative has roots
$\pm 97$ --- by $167$ gives $x^2(x-308)(x-360)$, the quartic of Guy's
first printed pair; normalizing to the form $x^2(x-1)(x-a)$ gives
$a=360/308=90/77$.  Guy's second printed pair $(668668,1990440)$ equals
$4\cdot(167167,497610)$, and its quartic --- first-derivative roots
$\{0,423696,1570635\}$, second-derivative roots $\{195624,1133930\}$,
all integers --- normalizes to $a=497610/167167$ (lowest terms).  With
$P=(-12,36)$ (equivalently $(-8,36)$ on $E_0$ under $w=x-4$), direct
computation of the multiples gives the four parameters displayed, with
both gates of \eqref{eq:gates} satisfied at each of $2P$, $3P$, $4P$,
$5P$ (script and log archived with the release accompanying this
paper).  Guy's third printed pair
$(-277132848044,\,514660109040)$ normalizes to $a(4P)$ exactly; his
fourth,
\[
(16695809521921862640,\ 28041466545675190604),
\]
normalizes to $1/a(5P)$, i.e.\ to $a(5P)$ after exchanging the two
simple roots.  The four printed quartics are thus, up to that symmetry,
the images of the four consecutive multiples $2P$ through $5P$.

The family's record from 1950 to 2019 --- Klamkin's isolated example,
Carroll's uniqueness conjecture \cite{Carroll1989}, the pairs printed
in Guy's column \cite{Guy1989} (with a refutation credited separately
to M.~J.~Miller, whose example we do not identify), the table of
\cite{BK1995} --- which already lists these four quartics as the
consecutive multiples $n\,(75,405)$, $n=1,\dots,4$, on $E_{A}$,
identifying the first with Carroll's example, the organizing
observation in print --- the rank tables of \cite{BM2000}, and the
infinite families of \cite{BC2016,Carrillo2019} --- may be viewed
retrospectively as the discovery, one multiple at a time, of the images
of a single point on a single curve.  The steps added here are the
identification with Guy's printed pairs, the transport to the single
curve $E_{0}$, and the identification of Stroeker's Jacobian
(Theorem~\ref{thm:ranks}).

%% ================================================================
%% PART II — THE REPAIR
%% ================================================================
\part*{Part II. The repair}

%% ----------------------------------------------------------------
%% §4 THEOREM 7′ (~5p)
%% SOURCE: RNT §2–4 全量 + blueprint Ch3
%% Dvořák 方式: 定理見出しに formal 宣言名併記
%% ----------------------------------------------------------------
\section{Theorem $7'$: the unconditional repair}\label{sec:repair}

Part I located the unproved step; this part states the classification
that the evidence supports and proves it unconditionally.  The 1995
kernel of the result --- the description, due to Buchholz and Kelly
\cite{BK1995}, of the family $x^2(x-1)(x-a)$ --- is sound, and the gap
lies precisely in the 2000 increment; Theorem~$7'$ below is its
constructive counterpart.  Two features of the proof are worth
announcing.  First, \emph{no rank computation enters}: the published
treatment proceeds via the Mordell--Weil rank of the curve, and we show
that the rank is needed only to \emph{enumerate} its rational points,
not to establish the classification itself.  Both directions reduce to
explicit polynomial certificates together with elementary field and
order arguments.  Second, the entire theorem is
formalized: the main statement
\leanline{Thm7Prime/Master.lean}{619}{thm7prime} compiles against
mathlib with the standard three-axiom footprint (\S\ref{sec:formal}).  The main text records each identity and the
logical role of each certificate; the archived development carries the
complete coefficient lists.

\subsection{The normalised family and the statement}
\label{sec:repair:statement}

{\sloppy
A quartic of type $p(2,1,1)$ with double root $c$ and simple roots
$r_1 \neq r_2$ is carried by the affine substitution
$x = c + (r_1 - c)u$ to a nonzero multiple of the \emph{normalised
family}
\begin{equation}\label{eq:family}
q_a(x) \;=\; x^{2}(x-1)(x-a), \qquad a = \tfrac{r_2-c}{r_1-c},
\end{equation}
with $a \notin \{0,1\}$; rational-derivedness is preserved by such
substitutions.  (This normalisation identity is itself one of the
formalized lemmas,
\leanline{Thm7Prime/Master.lean}{349}{caseA\_normalized}; the full
bridge from an arbitrary split $(2,1,1)$ quartic is the affine
classification of \S\ref{sec:fidelity}.)  Since
$q_a = x^4 - (1{+}a)x^3 + a x^2$, we have
\[
q_a' = x\,(4x^2 - 3(1{+}a)x + 2a), \qquad
q_a'' = 12x^2 - 6(1{+}a)x + 2a,
\]
and $q_a'''$ always splits.  Hence $q_a$ is rational-derived if and
only if the two discriminant \emph{gates}
\begin{equation}\label{eq:gates}
G_1(a):\; 9a^{2}-14a+9 = r^{2}, \qquad
G_2(a):\; 9a^{2}-6a+9 = s^{2}
\end{equation}
admit rational solutions $r, s$
(\leanline{Thm7Prime/Master.lean}{70}{Gate1},
\leanline{Thm7Prime/Master.lean}{73}{Gate2}).  Write $\mathrm{RD}(a)$
for the conjunction ``$a \neq 0$, $a \neq 1$, $G_1(a)$, $G_2(a)$''
(\leanline{Thm7Prime/Master.lean}{77}{RD211}); the equivalence of this
gate form with the natural derivative-splitting form is itself a kernel
theorem (\S\ref{sec:fidelity}).\par}

On the curve
\begin{equation}\label{eq:E}
E\colon\quad z^{2} \;=\; w^{3}+12w^{2}-108w \;=\; w(w-6)(w+18)
\qquad(\text{Cremona } \texttt{576i2}),
\end{equation}
the identity
$(z-w-18)(8w+z) = (7w-18)z + (w^3+4w^2-252w)$ holds
(\leanline{Thm7Prime/Master.lean}{95}{D\_reduce}), and we write
$D(w,z)$ for this common value of the denominator of the
Buchholz--MacDougall map
\begin{equation}\label{eq:amap}
a(w,z) \;=\; \frac{9\,(2w+z-12)(w+2)}{(z-w-18)(8w+z)}
\end{equation}
(\leanline{Thm7Prime/Master.lean}{160}{aMap}).  Here and below,
$(w,z) \in E(\Q)$ means an affine rational point of \eqref{eq:E}; the
point at infinity, at which \eqref{eq:amap} is undefined, plays no
role.

\begin{theorem}[Theorem $7'$, image form (the normalised $p(2,1,1)$
sector);
\leanline{Thm7Prime/Master.lean}{619}{thm7prime}]\label{thm:main}
Let $a \in \Q$.
\begin{enumerate}
\item[(i)] \textup{(Backward.)} If $(w,z) \in E(\Q)$ with
$D(w,z) \neq 0$ and $a := a(w,z) \notin \{0,1\}$, then $\mathrm{RD}(a)$.
\item[(ii)] \textup{(Forward / completeness.)} If $\mathrm{RD}(a)$, then
there exists $(w,z) \in E(\Q)$ with $D(w,z) \neq 0$ and $a(w,z) = a$.
\end{enumerate}
\end{theorem}

Interchanging the two simple roots of \eqref{eq:family} replaces $a$ by
$1/a$, so $q_a$ and $q_{1/a}$ represent the same affine-equivalence
class; Theorem~\ref{thm:main} characterizes the complete set of
normalised parameters, not a choice of representative for each class.
The theorem is stated on $E = \texttt{576i2}$ with \emph{no coset
restriction}; the relation to the coset-restricted 1995 statement is
\S\ref{sec:repair:bk}.  The Weierstrass model is tied to mathlib's own
vocabulary inside the development
(\leanline{Thm7Prime/Master.lean}{288}{OnE\_iff\_equation},
\leanline{Thm7Prime/Master.lean}{296}{E576i2\_elliptic}).

\subsection{The backward direction: gate identities}
\label{sec:repair:backward}

The backward direction is a pair of square identities in the function
field of $E$.  Set $M(w) = w(w+18)(w^2-63w+486)$; note
$w^2-63w+486 = (w-9)(w-54)$.  On $E$, away from the loci
$M(w) = 0$ and $D(w,z) = 0$, one has the \emph{exact} identities
\begin{align}
9a^2 - 14a + 9 &=
\Big(\tfrac{3(w-18)(w+6)w(w+18) + (-21w^2+108w-2916)\,z}{M(w)}\Big)^{2},
\label{eq:g1sq}\\
9a^2 - 6a + 9 &=
\Big(\tfrac{3(w-18)(w+6)w(w+18) + (-9w^2-648w+2916)\,z}{M(w)}\Big)^{2},
\label{eq:g2sq}
\end{align}
where $a = a(w,z)$.  Cleared of denominators these are polynomial
identities modulo the curve equation, verified in Lean by
\texttt{linear\_combination} with explicit cofactor certificates
(\leanline{Thm7Prime/Master.lean}{99}{gate1\_cleared},
\leanline{Thm7Prime/Master.lean}{108}{gate2\_cleared}); the
certificates were computed exactly in Sage and are reproduced verbatim
in the formal development, so their correctness is checked by the
kernel rather than trusted.

The excluded locus is finite and harmless: on $E(\Q)$ the conditions
$M(w) = 0 \neq D(w,z)$ hold only at $(54,432)$ and $(9,-27)$, both of
which map to $a = 77/90$, where the gates hold with the explicit
witnesses $r = 19/10$ and $s = 97/30$
(\leanline{Thm7Prime/Master.lean}{195}{exceptional\_points}).  This
proves Theorem~\ref{thm:main}(i)
(\leanline{Thm7Prime/Master.lean}{255}{thm7\_backward}).

\subsection{The forward direction: an explicit reconstruction from the
gate witnesses}\label{sec:repair:forward}

Let $a$ satisfy $\mathrm{RD}(a)$, with gate witnesses $r, s$ as in
\eqref{eq:gates}; note the sign freedom $r \mapsto -r$, $s \mapsto -s$.
Gr\"obner elimination shows that a point of $E$ above $a$ can be
reconstructed rationally from the triple $(a,r,s)$.  Concretely: in
$\Q[z,w,a,r,s]$ with lexicographic order, let $I$ be generated by the
curve equation, the graph relation $a\,D(w,z) - 9(2w+z-12)(w+2)$, and
the relations $r\,M(w) - N_1(w,z)$, $s\,M(w) - N_2(w,z)$ tying the
witnesses to the square roots of \eqref{eq:g1sq}--\eqref{eq:g2sq},
with $N_1, N_2$ the numerators displayed there.  After saturating $I$
by $D(w,z)\,M(w)$, its Gr\"obner basis contains the linear relation
$(3r-7s+12)\,w - 18\,(15a-6r-s-13)$.  The saturated Gr\"obner-basis
computation was performed in Magma; Sage independently verified the
resulting reconstruction formulas and the polynomial identities used
below.  The Gr\"obner computation is used only for \emph{discovery} ---
correctness of the reconstruction is established by the kernel-checked
identities.  Set
\begin{equation}\label{eq:winv}
W = 18\,(15a - 6r - s - 13), \qquad T = 3r - 7s + 12, \qquad
w \;=\; \frac{W}{T}
\end{equation}
(\leanline{Thm7Prime/Master.lean}{371}{wNp},
\leanline{Thm7Prime/Master.lean}{374}{wDp}),
and determine $z$ linearly from $a \cdot D(w,z) = 9(2w+z-12)(w+2)$:
explicitly $z = z_N/z_D$ with
\begin{equation}\label{eq:zN}
\begin{aligned}
z_N &= 9\,(2W-12T)(W+2T)\,T - a\,\bigl(W^{3}+4W^{2}T-252\,WT^{2}\bigr),\\
z_D &= 18\,K\,T^{2},
\end{aligned}
\end{equation}
where
\begin{equation}\label{eq:K}
K(a,r,s) \;=\; 105a^{2}-45ar-238a+51r+16s+105
\end{equation}
(\leanline{Thm7Prime/Master.lean}{382}{zNp},
\leanline{Thm7Prime/Master.lean}{387}{zDp},
\leanline{Thm7Prime/Master.lean}{378}{K1p}).  Two polynomial
consequences are then kernel-checked modulo the two gate relations,
each by a single \texttt{linear\_combination} step with a Sage-derived
cofactor certificate:
\leanline{Thm7Prime/Master.lean}{511}{curve\_cleared} proves that the
constructed point lies on $E$, and
\leanline{Thm7Prime/Master.lean}{525}{aden\_cleared} proves the
pullback identity
\begin{equation}\label{eq:pullback}
D(w,z)\cdot T^{3} z_D \;=\; -62208\; P(a,r,s)
\end{equation}
for the explicit sextic $P$ of Appendix~\ref{app:data}
(\leanline{Thm7Prime/Master.lean}{392}{adP6}), so that
$P(a,r,s) \neq 0$ together with $T \neq 0$ and $z_D \neq 0$ forces
$D(w,z) \neq 0$.  Independently, the definition of $z$ gives the
cleared graph identity $a\,D(w,z) = 9(2w+z-12)(w+2)$; once the sign
choice and \eqref{eq:pullback} establish $D(w,z) \neq 0$, this identity
yields $a(w,z) = a$.

Three denominators must be controlled, and this is where the sign
freedom is spent.
\begin{enumerate}
\item[(a)] \emph{The $w$-denominator.}  Write
$Q(a) = 225a^{2}-210a+238 = (15a-7)^{2}+189 > 0$.  An explicit
elimination certificate gives
$a^{2}Q(a) \in \langle\, T,\ r^{2}-(9a^{2}-14a+9),\
s^{2}-(9a^{2}-6a+9) \,\rangle$, so for $a \neq 0$ the gate relations
force $T \neq 0$, for \emph{every} choice of signs
(\leanline{Thm7Prime/Master.lean}{424}{wD\_ne}).
\item[(b)] \emph{The factor $K$.}  A second elimination certificate
gives
$(3a-5)(5a-3)\,a^{2}Q(a) \in \langle\, K,\
r^{2}-(9a^{2}-14a+9),\ s^{2}-(9a^{2}-6a+9) \,\rangle$.  At $a = 5/3$
the first gate forces $(3r)^{2} = 96$, and at $a = 3/5$ it forces
$(5r)^{2} = 96$; since $96$ is not a rational square, these branches
are vacuous, and again $K \neq 0$ for every choice of signs once
$a \neq 0$ (\leanline{Thm7Prime/Master.lean}{448}{K1\_ne},
\leanline{Thm7Prime/Master.lean}{405}{sq\_ne\_96}).
\item[(c)] \emph{The $a$-map denominator at the constructed point.}
By the pullback identity \eqref{eq:pullback}, this reduces to
$P(a,r,s) \neq 0$.  The ideal generated by the two gate relations
together with the four sign conjugates $P(a,\pm r,\pm s)$ is the
\emph{unit ideal}: a Nullstellensatz certificate (six cofactors,
reproduced in the formal development) exhibits $1$ as a combination,
so at least one sign choice keeps $P \neq 0$
(\leanline{Thm7Prime/Master.lean}{489}{adP6\_cover}).
\end{enumerate}
Choosing signs by (c) and applying (a), (b) yields a point of $E(\Q)$
off the degenerate locus with $a(w,z) = a$
(\leanline{Thm7Prime/Master.lean}{542}{forward\_core},
\leanline{Thm7Prime/Master.lean}{597}{thm7\_forward}), proving
Theorem~\ref{thm:main}(ii).  For example $a = 90/77$, with
$r = 171/77$, $s = 291/77$, yields
$(w,z) = (726/25,\, 22176/125)$.

\begin{remark}[No rank input]\label{rem:norank}
The published treatment of Theorem~7 proceeds through the
Mordell--Weil rank of $E$ (equal to $1$; re-verified unconditionally by
$2$-descent in two independent systems).  Theorem~$7'$ consumes no rank
fact in either direction: rank enters only if one wishes to
\emph{enumerate} $E(\Q)$ and hence list the rational-derived
parameters, as in \cite{BK1995}.  In the formal development the
rank-$1$ statement is present as a single disclosed axiom,
instance-localized to the curve
(\leanline{Thm7Prime/Master.lean}{320}{rank\_E576i2}), retained for
that enumerative purpose and consumed only by auxiliary declarations
from an earlier conditional architecture; the kernel's per-declaration
axiom report records that the main theorem does not depend on it
(\S\ref{sec:formal}).
\end{remark}

\subsection{Relation to the 1995 classification}\label{sec:repair:bk}

Theorem~$7'$ is not a transport of the theorem of \cite{BK1995} across
the isogeny of Remark~\ref{rem:isogeny}, for two reasons worth making
explicit.  First, the 1995 classification runs through the subgroup
structure of the model curve \texttt{576i3}: the set of points
corresponding to rational-derived quartics is the index-two subgroup
generated by $(75,405)$ and $(-6,0)$ (their Theorem~1), the parameter
map $(X,Y) \mapsto (5X+Y+30)/(9(X+2))$ is injective on the infinite
cyclic subgroup generated by $(75,405)$ (their Lemma~4), and the
passage between the two cosets is the substitution symmetry
$a \mapsto 1/a$ of the family (their Lemma~3).  Neither subgroup
carries an intrinsic description in the statement; both are identified
through the $2$-descent and generator computation of that paper, so the
classification, and any statement transported from it, consumes the
Mordell--Weil enumeration of the curve.  Theorem~$7'$ consumes no such
input: on \texttt{576i2} with the map \eqref{eq:amap}, no subgroup or
coset restriction is needed, and the degenerate locus is finite and
explicit.  Second, the completeness direction of \cite{BK1995} passes
through the quartic curve $Y^2 = (9X^2-14X+9)(9X^2-6X+9)$, on which
only the \emph{product} of the two gates is a square; deciding which of
its points satisfy each gate individually is again the subgroup
determination.  The reconstruction \eqref{eq:winv} works directly from
the pair of gate witnesses and is, to our knowledge, new.  What
Theorem~$7'$ shares with \cite{BK1995} is the normalised family and the
two gates; what it removes is the dependence of the classification on
the arithmetic of the curve.

The closest precursor of the rank-free route, however, is a sketch in
\cite{BM2000} itself.  Working over quadratic fields, \S3 of that
paper solves one gate by the chord method, substitutes the resulting
parametrization into the other, and passes through Mordell's
transformation to the same curve $E$ and to precisely the map
\eqref{eq:amap}, asserting that the map provides ``the correspondence
between all $k$-rational points \dots\ and all $k$-derived quartics''
before rank determination enters --- for enumeration.  Read literally,
the asserted correspondence does not hold: the map is undefined where
$\mathrm{aden}$ vanishes, its values $0$ and $1$ correspond to
degenerate quartics (the $\Q$-fibres are computed in
\S\ref{sec:fidelity:excluded}), and neither direction of the correspondence
is argued there.  Theorem~$7'$ can be read as the proved form of that
sentence: both directions are established, the exceptional locus is
exactly characterized, and the certificates are closed-form, hence
kernel-checkable.  The comparison also sharpens \S\ref{sec:ranks}:
that the rank should be needed only to \emph{enumerate} is implicit in
the presentation of \cite{BM2000} itself; what was missing was the
proof.

\begin{remark}[The two curves]\label{rem:isogeny}
The Buchholz--Kelly parametrization lives on
$E_A = \texttt{576i3}\colon Y^2 = X(X-48)(X+6)$ and is stated there
with a two-coset restriction; the map \eqref{eq:amap} lives on
$E = \texttt{576i2}$ and, by Theorem~$7'$, needs no restriction.  The
two are reconciled by the explicit $2$-isogeny $E \to E_A$,
\[
(w,z) \;\longmapsto\;
\Bigl( \frac{(w+6)^{2}}{w-6},\;
\frac{z\,(w+6)(w-18)}{(w-6)^{2}} \Bigr),
\]
with kernel $\{O, (6,0)\}$, verified independently in Magma and Sage.
The coset restriction appearing in \cite{BK1995} is a phenomenon of
their model and their parameter map, not of
\eqref{eq:E}--\eqref{eq:amap}.
\end{remark}

\begin{remark}[Errata in the sources]\label{rem:errata}
For the convenience of readers of \cite{BK1995} we record, verified by
direct recomputation in the computer-algebra layer (script and log
archived with the release accompanying this paper): rows $n=2,3$ of the table of \cite{BK1995}
(p.~131) are internally inconsistent --- the printed \emph{points} and
the printed $a$-values correspond to \emph{opposite} multiples $\pm nQ$
of the printed generator $Q=(75,405)$: the printed points equal $-2Q$
and $-3Q$, while the printed $a$-values and root pairs correspond to
$+2Q$ and $+3Q$; no identification $(X,Y)\equiv(X,-Y)$ that would
reconcile them is stated in that paper.  The misprinted constant
of \cite{BM2000} is \S\ref{sec:gap:provenance}; its two corrected
Mordell--Weil ranks are \S\ref{sec:ranks}.  None of this affects the
mathematics above.
\end{remark}

%% ----------------------------------------------------------------
%% §5 FIDELITY, CLASSIFICATION, ELEMENTARY SECTORS (~3p)
%% SOURCE: RNT Remark 3.2 + Fidelity/Classification.lean + jnt §2.3–2.4 相当
%% ----------------------------------------------------------------
\section{Statement fidelity and the affine classification}
\label{sec:fidelity}

The kernel guarantees that a formal proposition follows from the
axioms; whether that proposition models the mathematical problem is
outside the kernel.  This correspondence --- what the AFM's editorial
guidance calls \emph{proof auditing} --- was treated in this project as
a gate: every statement-level definition was cross-checked in Sage and Magma against
independent reimplementations \emph{before} any proof was attempted,
and the two links most exposed to a statement-level slip were
subsequently closed in the kernel as well.  This section records them,
together with the elementary sectors and the sanity controls.

\subsection{Gate form and natural form}\label{sec:fidelity:gate}

The gate form $\mathrm{RD}(a)$ of \S\ref{sec:repair:statement} was
chosen at design time as the division-free, verification-friendly
restatement of rational-derivedness.  Its correspondence with the
natural form --- the polynomial $Q_a = X^{2}(X-1)(X-\mathrm{C}\,a)$ and
its first three \texttt{Polynomial.derivative}s split over $\Q$
(\leanline{Thm7Prime/Fidelity.lean}{27}{NaturalRD}) --- is a kernel
theorem:

\begin{theorem}[Gate form $\iff$ natural form;
\leanline{Thm7Prime/Fidelity.lean}{173}{rd211\_iff\_natural}]
\label{thm:fidelity}
$\mathrm{RD211}(a) \iff a\neq 0 \wedge a\neq 1 \wedge
\mathrm{NaturalRD}(a)$.
\end{theorem}

Forward, the kernel-computed derivatives split explicitly under the
gates; conversely, a quadratic over $\Q$ with nonzero leading
coefficient that splits has square discriminant, applied to both
derivatives.  The CAS statement-fidelity gate remains in the
development as the independent design-time cross-check
(\texttt{scripts/sage/wp1\_fidelity.sage}).

\subsection{The affine classification}\label{sec:fidelity:affine}

The reviewer-facing form of the repaired result is the bridge from an
arbitrary split $(2,1,1)$ quartic to the normalised family, packaged as
one kernel theorem.  With $f\sim g$ iff $g=\nu\,f(\lambda X+\mu)$ for
$\lambda,\nu\neq 0$ --- an equivalence relation, kernel-checked
(\leanline{Thm7Prime/Classification.lean}{30}{AffineEquiv}) --- and
$f$ a split $(2,1,1)$ quartic
(\leanline{Thm7Prime/Classification.lean}{35}{IsP211}):

\begin{theorem}[Affine classification;
\leanline{Thm7Prime/Classification.lean}{175}{classification},
\leanline{Thm7Prime/Classification.lean}{188}{classification\_by\_curve}]
\label{thm:classification}
For a split $(2,1,1)$ quartic $f$: $f$ is rational-derived iff
$f\sim Q_{a}$ for some $a$ with
$\mathrm{RD211}(a)$; equivalently, iff $f\sim Q_{a(w,z)}$ for an affine
rational point $(w,z)$ of the Weierstrass model of \texttt{576i2} with
$D(w,z)\neq0$ and $a(w,z)\notin\{0,1\}$.
\end{theorem}

Rational-derivedness is invariant under $\sim$ (chain rule and
transport of splitting along linear substitutions; $\sim$ is the
equivalence relation induced by the action of the group
$\langle X\rangle$ of \cite{BM2000} --- the conservative vertical
translations extending it to $\langle X^{*}\rangle$ are exactly the
subject of \S\ref{sec:gap}); the
normalisation $a=(r_{2}-c)/(r_{1}-c)$ is the kernel form of
\texttt{caseA\_normalized}; the gates then come from
Theorem~\ref{thm:fidelity}, and the curve form from
Theorem~\ref{thm:main}.

\subsection{The excluded parameters and the elementary sectors}
\label{sec:fidelity:excluded}

The exclusion $a \neq 0$ is genuinely needed and completely understood:
the fibre $a^{-1}(0) \cap E(\Q)$ consists of the three points $(6,0)$,
$(-12,36)$, $(-2,-16)$, corresponding to the degeneration of the family
to the (always rational-derived) type $p(3,1)$.  The exclusion
$a \neq 1$ is vacuous on rational points: an exact divisor computation
shows the fibre $a^{-1}(1)$ has \emph{no} rational points (the two
quadratic factors of its eliminant $(w+2)^2(w^2-36w-108)(w^2-108)$ have
no rational roots, and at $w=-2$ the graph relation forces $z=16$,
where $D(-2,16)=0$, outside the domain of the map); independently,
$a=1$ fails the second gate since $12$ is not a square.  The elementary
sectors are dispatched directly: types $p(4)$ and $p(3,1)$ are always
rational-derived (up to the substitutions of
\S\ref{sec:repair:statement} these are $x^{4}$ and $x^{3}(x-b)$, whose
derivatives have visibly rational roots), and type $p(2,2)$ never is
(for $x^2(x-t)^2$ one finds $\operatorname{disc}(q'') = 48t^{2}$, and
$48$ is not a square).  The pieces assemble into the full statement:

\begin{theorem}[Theorem $7'$, assembled form: the repeated-root
classification]\label{thm:assembled}
A quartic over $\Q$ with a repeated root is rational-derived exactly in
the following cases: types $p(4)$ and $p(3,1)$, always; type $p(2,2)$,
never; and, for a split $(2,1,1)$ quartic, if and only if it is
$\sim$-equivalent to $Q_{a(w,z)}$ for an affine rational point $(w,z)$
of \texttt{576i2} with $D(w,z)\neq0$ and $a(w,z)\notin\{0,1\}$.
\end{theorem}

\noindent Each clause decomposes into kernel theorems: the $p(2,1,1)$
clause is Theorem~\ref{thm:classification}, through
Theorem~\ref{thm:fidelity} and the image form Theorem~\ref{thm:main};
the elementary sectors are the computations above; the case split by
multiplicity type, and the observation that a quartic with an
irrational root is not rational-derived (so the split hypothesis is
harmless), are elementary [M].  The remaining sector --- four distinct
roots --- is Part~III.

\subsection{Counterexamples and sanity controls}
\label{sec:fidelity:sanity}

Following the practice of exhibiting the counterexamples that delimit a
formalized statement, the development carries a sanity
layer, kernel-checked independently of the main proofs: a positive
control (the parameter $a=77/90$ passes both gates,
\leanline{Thm7Prime/Master.lean}{251}{gates\_at\_77\_90}) and a
near-miss that doubles as the negative control (Stroeker's properly
$K$-derived quartic of Theorem~\ref{thm:gap}: its $f'$ splits over
$\Q$ while $f''$ does not,
\leanline{Gap/BMThm7GapK.lean}{142}{fp\_splits},
\leanline{Gap/BMThm7GapK.lean}{148}{fpp\_not\_split\_over\_Q}, so the
quartic passes the first gate over $\Q$, fails the second, and the pair
of gates --- not either alone --- carries the classification).  The same quartic serves as the transparent
negative control of the transcript (\S\ref{sec:gap}), read over two
base fields; the sanity layer and the audit share their witnesses by
design, so that a statement-level slip in either would have been caught
by the other.

%% ================================================================
%% PART III — THE BOUNDARY
%% ================================================================
\part*{Part III. The boundary}

%% ----------------------------------------------------------------
%% §6 TERRAIN AND THE NORMALIZATION BRIDGE (~5–6p)
%% SOURCE: blueprint Ch4 (v1.0.4) — 用語は repo と一致 (査読者は突合する)
%% [K]/[C]/[M]/[A] ラベル規約をこの部の冒頭で宣言 (Ch4 と同文)
%% ----------------------------------------------------------------
\section{The terrain and the normalisation bridge}\label{sec:boundary}

With Theorem~$7'$ the repeated-root sector is classified
unconditionally; what remains open of the original Theorem~7 is exactly
its $p(1,1,1,1)$ clause, entangled with Conjecture~1 of \cite{BM2000}.
This part maps the terrain surrounding that boundary.  Throughout it,
each load-bearing claim carries a label: [K] = kernel-verified in Lean;
[C] = computer-algebra computation (two-system where so stated; scripts
and coverage in \texttt{scripts/CERTIFICATES.md}); [M] = ordinary
mathematical deduction, written but not kernel-checked; [A] = research
assessment.  Narrative sentences without a label carry no evidential
weight.

\subsection{The chart, the splitting surface, and the three covers}
\label{sec:boundary:chart}

Fixed notation for the whole part.  $B$ denotes the rational
$(a,b)$-chart of the $\sigma_{3}=0$ root space of \S2.2.3 of
\cite{BM2000}: the quartic has roots $\{1,a,b,c\}$ with
$c=-ab/\mathrm{den}$, $\mathrm{den}=ab+a+b$; its function field is
$K(B)=\Q(a,b)$.  The chart quantities
(\leanline{Gap/BMThm7Boundary.lean}{19}{den},
\leanline{Gap/BMThm7Boundary.lean}{21}{s1n},
\leanline{Gap/BMThm7Boundary.lean}{25}{Rq},
\leanline{Gap/BMThm7Boundary.lean}{27}{Sq},
\leanline{Gap/BMThm7Boundary.lean}{29}{Q4}) are
$s_{1n}=\sigma_{1}\,\mathrm{den}$, $s_{2n}=\sigma_{2}\,\mathrm{den}$,
the cleared splitting conditions
$R_{q}=\mathrm{den}^{2}(9\sigma_{1}^{2}-32\sigma_{2})$ (for $f'$) and
$S_{q}=\mathrm{den}^{2}(9\sigma_{1}^{2}-24\sigma_{2})$ (for $f''$), and
the funnel
$Q_{4}=(3R_{q}-S_{q})/18=\mathrm{den}^{2}(\sigma_{1}^{2}-4\sigma_{2})$.
Three kernel identities organize everything below:
\begin{align}
18\,Q_{4}&=3R_{q}-S_{q}
&&\text{(I1, \leanline{Gap/BMThm7Boundary.lean}{36}{I1\_cleared})},
\label{eq:I1}\\
s_{1n}^{2}-R_{q}&=-8\,Q_{4}
&&\text{(I2, \leanline{Gap/BMThm7Boundary.lean}{40}{I2})},
\label{eq:I2}\\
Q_{4}&=F_{1}\,F_{2}\,F_{3}
&&\text{(\leanline{Gap/BMThm7Boundary.lean}{44}{Q4\_factor})},
\label{eq:Q4factor}
\end{align}
with $F_{1}=ab+a+b-1$, $F_{2}=ab+a-b^{2}+b$, $F_{3}=a^{2}-ab-a-b$.
The kernel statements are the identities only; that $Q_{4}$ is
squarefree and that each $F_{i}=0$ is an irreducible plane curve of
genus~$0$ are two-system computer-algebra facts [C], not part of the
kernel statement.  Three further kernel factorisations
(\leanline{Gap/BMThm7Boundary.lean}{50}{pairing1},
\leanline{Gap/BMThm7Boundary.lean}{55}{pairing2},
\leanline{Gap/BMThm7Boundary.lean}{60}{pairing3}) identify, away from
the vanishing of explicit cofactors, the locus $Q_{4}=0$ --- where
Theorem~7's conclusion is algebraically forced --- with the locus on which the
quartic $\{1,a,b,c\}$ is symmetric under one of the three root
pairings: Buchholz--MacDougall's own excluded case.

$S$ denotes the \emph{splitting cover} of $B$: the biquadratic cover
obtained by adjoining $z$ with $z^{2}=R_{q}$ and $w$ with
$w^{2}=S_{q}$; its function field is, by definition,
$K(S)=\Q(a,b)(z,w)$.  Concretely, let $U_{B}\subset B$ be the complement of the vanishing
locus of
\[
\mathrm{den}\cdot ab\,(a-1)(b-1)(a-b)\,(ab+\mathrm{den})\,
(a+\mathrm{den})\,(b+\mathrm{den})
\]
--- exactly the base-variable inequalities of a nondegenerate terrain
point (Proposition~\ref{prop:bridge}: the factors beyond
$\mathrm{den}$ express, after clearing $\mathrm{den}$, that the four
roots $\{1,a,b,c\}$, $c=-ab/\mathrm{den}$, are pairwise distinct and
nonzero).  The splitting surface over the chart is the affine model
\[
S_{U}:=\operatorname{Spec}\bigl(\mathcal O(U_{B})[z,w]/
(z^{2}-R_{q},\,w^{2}-S_{q})\bigr),
\]
and we write $U:=S_{U}$ for this open model.  Every element of
$K(S)$ is a $K(B)$-combination of the spanning set $\{1,z,w,zw\}$ [M];
it is a basis once Proposition~\ref{prop:S_integral} is in place.
Kernel statements about $K(S)$ below are proved in span form ---
quantifying over \emph{any} field equipped with an embedding of $K(B)$
and square roots of $R_{q}$, $S_{q}$ --- and reach all of $K(S)$
through that spanning set [M].

For a critical point $x_{0}$ of the chart quartic $q$, \emph{the
disjunct at $x_{0}$} is the statement that the vertical translate
$q-q(x_{0})$ --- the transcript's \texttt{bmTranslate}, the very
operation the printed proof performs --- splits over $\Q$.  The
critical points are $0$ and $(3\sigma_{1}\pm\rho)/8$ with
$\rho=z/\mathrm{den}$; the corresponding disjuncts are $D_{0}$,
$D_{+}$, $D_{-}$
(\leanline{Gap/BMThm7Fibre.lean}{39}{D0},
\leanline{Gap/BMThm7Fibre.lean}{43}{Dplus},
\leanline{Gap/BMThm7Fibre.lean}{47}{Dminus}).  Exchanging the sign of
the square root $\rho$ exchanges $D_{+}$ and $D_{-}$; no conditions
beyond $\mathrm{den}\neq0$ and the splitting witness $z^{2}=R_{q}$ are
folded into the predicates.  Each disjunct is governed by a double
cover of $U$:

\begin{center}\small
\begin{tabular}{llll}
\toprule
disjunct & cover & radicand & exclusion \\
\midrule
$x_{0}=0$ & $\widetilde C_{0}\colon v^{2}=Q_{4}$ & $Q_{4}$ (by I1) &
[K] Thm~\ref{thm:multiquad} \\
$x_{0}=\tfrac{3\sigma_{1}+\rho}{8}$ & $C_{+}\colon v^{2}=\gamma_{+}$ &
$\gamma_{+}=s_{1n}(s_{1n}-z)$ & [K] Thm~\ref{thm:gamma} \\
$x_{0}=\tfrac{3\sigma_{1}-\rho}{8}$ & $C_{-}\colon v^{2}=\gamma_{-}$ &
$\gamma_{-}=s_{1n}(s_{1n}+z)$ & [K] Thm~\ref{thm:gamma} \\
\bottomrule
\end{tabular}
\end{center}

\noindent The two $C_{\pm}$ radicands are \emph{conjugate}; neither is
claimed to have square class $-2Q_{4}$ itself.  Their product ---
equivalently their layer norm --- has square class $-2Q_{4}$:
$\gamma_{+}\gamma_{-}=s_{1n}^{2}(s_{1n}^{2}-z^{2})=-8s_{1n}^{2}Q_{4}$
by \eqref{eq:I2}, and it is this \emph{norm class} that the kernel
obstruction consumes.  Three covers, one sextic, two obstruction
routes.  Each cover's branch locus is supported on the zero locus of
its radicand on $U$, with ramification exactly at the codimension-one
points where the radicand has odd valuation.

\subsection{The normalisation bridge}\label{sec:boundary:bridge}

In normalised form, Conjecture~1 asserts that no quartic
$x(x-1)(x-a)(x-b)$ with four distinct roots is $\Q$-derived
(\leanline{Gap/BMThm7Transcript.lean}{996}{Conjecture1\_normal};
stated as a definition, without proof --- the open statement of
\S\ref{sec:open}).  The chart of \S\ref{sec:boundary:chart} is where
the printed argument of \S2.2.3--2.2.4 of \cite{BM2000} works; the
following kernel theorem identifies the two pictures exactly.

\begin{proposition}[Normalization bridge;
\leanline{Gap/BMThm7Terrain.lean}{394}{conjecture1\_iff\_terrain}]
\label{prop:bridge}
There exists a $\Q$-derived quartic with four distinct rational roots
(the $\{0,1,a,b\}$ normal form of Conjecture~1) if and only if there
exists a nondegenerate rational point of the splitting terrain:
$(a,b,z,w)\in\Q^{4}$ with $z^{2}=R_{q}(a,b)$, $w^{2}=S_{q}(a,b)$,
$\mathrm{den}(a,b)\neq0$, and the four roots $\{1,a,b,c\}$,
$c=-ab/\mathrm{den}$, pairwise distinct and nonzero.  Equivalently:
the normalised Conjecture~1 holds iff $S(\Q)$ has no nondegenerate
point.  (The reduction of the unnormalised conjecture to the
$\{0,1,a,b\}$ normal form is the affine action --- ordinary
mathematics [M]; only the normal form enters the kernel.)
\end{proposition}

Both directions are kernel theorems, with every normalisation
condition tracked
(\leanline{Gap/BMThm7Terrain.lean}{124}{kderived\_of\_terrain},
\leanline{Gap/BMThm7Terrain.lean}{223}{terrain\_of\_kderived}).  In
outline: a terrain point yields the monic quartic
$f=(x-1)(x-a)(x-b)(x-c)$ with $e_{3}=0$, so $f'(0)=0$ and
$f'=x(4x^{2}-3e_{1}x+2e_{2})$; the witnesses $z,w$ split the quadratic
factors of $f'$ and $f''$ via
$9e_{1}^{2}-32e_{2}=(z/\mathrm{den})^{2}$ and
$9e_{1}^{2}-24e_{2}=(w/\mathrm{den})^{2}$, and an affine change carries
$f$ to the $\{0,1,a',b'\}$ normal form.  Conversely, a $\Q$-derived
quartic in normal form has a rational critical point $x_{0}$, which is
not a root (the roots are distinct, so $f$ is squarefree and
$f'(r)\neq0$ at each root $r$); translating $x_{0}$ to the origin and
scaling the nonzero root $-x_{0}$ to $1$ produces roots $\{1,A,B,C\}$
with $e_{3}=0$, forces $\mathrm{den}(A,B)\neq0$ (else $AB=0$,
contradicting nonzero roots), and the rational root differences of the
split quadratic factors supply the witnesses
$z:=4\,\mathrm{den}\,(r_{+}-r_{-})$,
$w:=6\,\mathrm{den}\,(u_{+}-u_{-})$.

\subsection{The fibre correspondence and the radicand derivation}
\label{sec:boundary:fibre}

The disjuncts and the covers are tied together fibrewise, in the
kernel.

\begin{proposition}[Fibre correspondence;
\leanline{Gap/BMThm7Fibre.lean}{293}{fibre\_correspondence}]
\label{prop:fibre}
Over $\mathrm{den}\neq0$ with $z^{2}=R_{q}$: $D_{0}$ holds iff the
fibre of the cleared model $\widetilde C_{0}\colon v^{2}=Q_{4}$ has a
rational point, and $D_{\pm}$ holds iff the fibre of
$C_{\pm}\colon v^{2}=\gamma_{\pm}$ has a rational point.  (The kernel
hypotheses are exactly these two --- $\mathrm{den}\neq0$ and
$z^{2}=R_{q}$; no further nondegeneracy is needed, so the formal
statement is slightly stronger than the chart discussion requires.)
\end{proposition}

For $D_{0}$
(\leanline{Gap/BMThm7Fibre.lean}{51}{D0\_iff\_disc}): since
$\sigma_{3}=0$, the translate at $0$ factors as
$q-q(0)=X^{2}(X^{2}-\sigma_{1}X+\sigma_{2})$, so $D_{0}$ holds iff
$\sigma_{1}^{2}-4\sigma_{2}$ is a rational square --- the fibre of the
\emph{chart} cover $C_{0}\colon t^{2}=\sigma_{1}^{2}-4\sigma_{2}$.  The
chart cover and the cleared model are then identified by
\leanline{Gap/BMThm7Fibre.lean}{98}{cleared\_model\_iff}: over
$\mathrm{den}\neq0$ the substitution $v=\mathrm{den}\,t$ is a bijection
on rational solutions (multiplication by the nonzero square
$\mathrm{den}^{2}$; the identity \eqref{eq:I1} is consumed
definitionally), so clearing the denominator introduces no spurious
solutions and loses none.  For $D_{\pm}$
(\leanline{Gap/BMThm7Fibre.lean}{117}{Dplus\_iff},
\leanline{Gap/BMThm7Fibre.lean}{204}{Dminus\_iff}) the radicand is
\emph{derived}, not posited: at
$x_{0}=(3s_{1n}\pm z)/(8\,\mathrm{den})$ the translate factors as
$(X-x_{0})^{2}(X^{2}+BX+C)$ with $B=2x_{0}-\sigma_{1}$,
$C=\sigma_{2}+3x_{0}^{2}-2\sigma_{1}x_{0}$, the linear-coefficient
match being exactly $q'(x_{0})=0$ (available from $z^{2}=R_{q}$); the
discriminant identity $B^{2}-4C=\sigma_{1}(\sigma_{1}\mp\rho)/4$ and
the clearing
$4\,\mathrm{den}^{2}(B^{2}-4C)=s_{1n}(s_{1n}\mp z)=\gamma_{\pm}$
(solution bijection $v=2\,\mathrm{den}\,d$) convert the splitting of
the quadratic cofactor into ``$\gamma_{\pm}$ is a rational square'' ---
the fibre of $C_{\pm}$.  Signs, denominators and the direction of
clearing are all carried inside the kernel statements; the same
identities are certified independently in CAS [C:
\texttt{terrain\_crosscheck}, seven checks covering the critical
points, the factorizations, the discriminant identities and the
clearing].

\subsection{Integrality of the splitting surface}
\label{sec:boundary:integral}

\begin{proposition}[Integrality]\label{prop:S_integral}
$R_{q}$, $S_{q}$ and $R_{q}S_{q}$ are nonsquares in $\Q(a,b)$, and
remain nonsquares in $\overline{\Q}(a,b)$.  Consequently
$[K(S):K(B)]=4$, the spanning set $\{1,z,w,zw\}$ is a basis, and the
splitting surface $S_{U}$ is geometrically integral, and its induced
dominant map to $B$ has generic degree four.
\end{proposition}

Over $\Q$ the three exclusions are kernel theorems
(\leanline{Gap/BMThm7NormCriterion.lean}{60}{Rq\_not\_square},
\leanline{Gap/BMThm7NormCriterion.lean}{67}{Sq\_not\_square},
\leanline{Gap/BMThm7NormCriterion.lean}{76}{RqSq\_not\_square}), by
value specialisation at witness points
($R_{q}(2,3)=2480$, $S_{q}(1,2)=876$, $(R_{q}S_{q})(1,2)=171696$, none
a rational square) [K].  A remark on method: on the chart, $R_{q}$ and
$S_{q}$ are positive semi-definite --- a quartic with four real roots
has non-negative derivative discriminants --- so any technique that
excludes squares by exhibiting a \emph{negative value} is unavailable
for them in principle, and value specialisation at non-square values
was the route adopted here.  Over $\overline{\Q}$: each of the three has a
nonconstant irreducible factor of odd multiplicity ($R_{q}$
irreducible; $S_{q}=3\cdot(\text{irreducible})$; the two are coprime)
[C], and a squarefree-base-change lemma [M] --- an irreducible factor
of odd multiplicity survives every algebraic extension of the constant
field, because squarefreeness is detected by a gcd computation
unchanged by extension, distinct factors share no geometric component,
and a square has even valuation along every divisor; an infinite
extension reduces to a finite subextension containing the finitely
many coefficients of a hypothetical square root --- shows none is a
square in $\overline{\Q}(a,b)$.  Degree four, the basis, and
integrality follow from the standard biquadratic criterion [M]: for
$r,s,rs$ all nonsquares in a field of characteristic $\neq2$, the
algebra $K[z,w]/(z^{2}-r,w^{2}-s)$ is a field of degree four with
basis $\{1,z,w,zw\}$.  Scheme-theoretically, the coordinate ring of
$S_{U}$ is a free $\mathcal O(U_{B})$-module on that basis, hence
torsion-free, hence embeds in its generic localisation --- the
degree-four field just constructed --- and a ring embedding into a
field is a domain; the same chain over $\overline{\Q}$ gives geometric
integrality [M].

\subsection{Cover nontriviality}\label{sec:boundary:covers}

A square in the function field $\Q(a,b)$ is non-negative at every
rational point where it is defined.  Each of the eight candidates whose
squareness would trivialise a double cover is \emph{negative} at an
explicit rational point --- the eight negativity certificates are
kernel lemmas at the points $(-51/13,-17/7)$ (a $K$-witness from
Theorem~\ref{thm:gap}) and $(1,3)$
(\leanline{Gap/BMThm7Boundary.lean}{72}{Q4\_neg} through
\leanline{Gap/BMThm7Boundary.lean}{94}{n2Q4RqSq\_neg}) --- and the
deduction itself is in the kernel: the specialisation principle is
evaluation along a ring homomorphism
(\leanline{Gap/BMThm7FunctionField.lean}{108}{sq\_eval\_nonneg}),
squares descend from the fraction field to the integrally closed ring
$\Q[a][b]$
(\leanline{Gap/BMThm7FunctionField.lean}{117}{square\_descent}), and
all eight non-squareness verdicts in $\Q(a,b)$ are kernel theorems
(\leanline{Gap/BMThm7FunctionField.lean}{134}{Q4\_not\_square} and its
seven companions).  The two-system record remains as the discovery
artifact [C].

For the two abstract kernel statements below, let $M$ be any field
containing $\Q(a,b)$ via an embedding
$\iota\colon\Q(a,b)\hookrightarrow M$, and let
$\varrho,\varsigma\in M$ satisfy $\varrho^{2}=\iota(R_{q})$ and
$\varsigma^{2}=\iota(S_{q})$ (written $\varrho,\varsigma$ to keep
them distinct from the chart quantity $\rho=z/\mathrm{den}$); all
spans are over $F:=\iota(\Q(a,b))$, with which we identify
$\Q(a,b)$ below --- in the kernel statements the four coefficients are
quantified over the base field.

\begin{theorem}[Non-squareness in the multiquadratic extension;
\leanline{Gap/BMThm7SquareClass.lean}{293}{Q4\_not\_square\_multiquadratic};
\leanline{Gap/BMThm7SquareClass.lean}{315}{n2Q4\_not\_square\_multiquadratic}]
\label{thm:multiquad}
In $M$ as above, neither $Q_{4}$ nor $-2Q_{4}$ is the square of any
element of
$\operatorname{Span}_{F}\{1,\varrho,\varsigma,\varrho\varsigma\}$.
\end{theorem}

The engine is a biquadratic square-class transfer
(\leanline{Gap/BMThm7SquareClass.lean}{81}{biquadratic\_transfer}):
if an element of the $F$-span squares to $x \in F$, then one of $x$,
$xR_{q}$, $xS_{q}$, $xR_{q}S_{q}$ is a square in $F$; degenerate cases are not excluded ---
the case analysis absorbs them.  The transfer reduces the theorem to
the eight base-field exclusions above.  Since every element of $K(S)$
lies in the span [M], the theorem excludes squares in all of $K(S)$;
the kernel boundary is exactly the span statement, the spanning fact
is ordinary mathematics.

\begin{theorem}[$C_{\pm}$-exclusion in the layer;
\leanline{Gap/BMThm7NormCriterion.lean}{272}{gamma\_not\_square\_multiquadratic}]
\label{thm:gamma}
For either sign, $\gamma_{\pm}=s_{1n}(s_{1n}\mp\varrho)$ --- an
element of the layer $L:=F(\varrho)$, not of the base
field --- is not, in $M$ as above, the square of any element of
$\operatorname{Span}_{F}\{1,\varrho,\varsigma,\varrho\varsigma\}$.
\end{theorem}

Because $\gamma_{\pm}$ lives in the layer, the biquadratic transfer
does not apply to it directly.  Instead a layer quadratic step
(\leanline{Gap/BMThm7NormCriterion.lean}{87}{step\_over\_layer})
reduces a hypothetical square root to a square in the layer or its
$S_{q}$-multiple, and the coordinate norm passage
(\leanline{Gap/BMThm7NormCriterion.lean}{236}{layer\_square\_norm}:
$g_{0}^{2}-g_{1}^{2}R_{q}=(p^{2}-q^{2}R_{q})^{2}$, computed in
coordinates, with no Galois action) turns either case, via
\eqref{eq:I2} in its polynomial form
(\leanline{Gap/BMThm7NormCriterion.lean}{38}{I2\_R2}), into the
statement that $-2Q_{4}$ is a square in $\Q(a,b)$ --- contradicting
the eight verdicts.  Both signs are handled at once.

Geometric integrality of the three covers follows [M with C and K
inputs]: $Q_{4}=F_{1}F_{2}F_{3}$ with distinct irreducible factors,
each to the first power, and $\gcd(F_{i},R_{q}S_{q})=1$ [C], so
$Q_{4}$ has odd valuation along every geometric component of each
$F_{i}=0$; the generic point of every such component lies outside the
branch divisor of $S\to B$ (equivalently, the two divisors share no
codimension-one component), and the valuations stay odd in
$\overline{\Q}(S)$; hence
$Q_{4}\notin\overline{\Q}(S)^{\times2}$, and the same divisor argument
applies verbatim to $-2Q_{4}$.  For $C_{\pm}$ the specialised kernel
theorem quantifies its coefficients over $\Q(a,b)$ and is \emph{not}
base-changed; instead the abstract kernel lemmas
(\texttt{step\_over\_layer}, \texttt{layer\_square\_norm} --- stated
for an arbitrary characteristic-zero field) are applied afresh with
$K=\overline{\Q}(a,b)$, their inputs supplied by the odd-multiplicity
divisor data through the base-change lemma, and the assembly reduces
geometric squareness of $\gamma_{\pm}$ to that of $-2Q_{4}$, excluded
above [M].  So $C_{0}$, $C_{+}$ and $C_{-}$ are dominant,
geometrically integral degree-two covers of $U$; removing the branch
locus yields \'etaleness, not integrality.  ($-2Q_{4}$ is not itself a
radicand: it enters as the layer-norm obstruction class of
$\gamma_{\pm}$.)

\subsection{What the terrain shows --- and what is not claimed}
\label{sec:boundary:reading}

The locus inside $S$ where Theorem~7's conclusion is algebraically
forced is the proper curve $S\cap\{Q_{4}=0\}$: six one-dimensional
components, each of genus~$0$ [C].  On each component $R_{q}$ restricts
to a square while $S_{q}$ does not; the restriction of $S_{q}$ has
constant square class $12=4\cdot 3$ [C], so $w^{2}=S_{q}$ has rational
solutions only where $S_{q}=0$, and every rational zero of $S_{q}$ on
each component is a degenerate point of the chart (on $F_{1}$ there are
no rational zeros at all) [C], hence inadmissible.  Hence over $\Q$ the
conclusion is never algebraically forced: \emph{for Theorem~7's conclusion to hold
at a nondegenerate $P\in S(\Q)$}, $P$ must lie in the image of at least
one of the three nontrivial double covers --- whether every terrain
point does so is precisely the open content of the theorem, and is not
assumed here.  Thus the algebraically forced locus $S\cap\{Q_{4}=0\}$ contributes
no admissible rational terrain point [C+M].

Away from the degeneracy locus each disjunct demands lifting through a
double cover that is nontrivial over the splitting extension $K(S)$, so
the image of the rational points of each cover is a thin subset of
type~II [M] --- essentially by definition, as the image of the rational
points of a dominant degree-two cover.  \emph{No assertion is made that
$S(\Q)$ contains rational points outside the union of these images}:
no such conclusion follows from thinness alone, and the additional
input that would provide it --- for example an appropriate
Hilbert-property statement for $S$ --- is not proved here; what is
proved is geometric integrality.  Thinness bounds what a generic-point argument can see,
not what rational points exist.  Within the approaches surveyed we
found no finite descent reduction, and make no claim that none exists;
and the field-dependence begins
already at Buchholz--MacDougall's own excluded case, over the quadratic
extension suggested by the square class, $\Q(\sqrt{3})$
($12=4\cdot 3$), while Stroeker's points live over $\Q(\sqrt{3441})$.

The reading is an assessment [A] anchored by machine artifacts, not a
theorem: the computations above eliminate the explicitly listed
algebraically forced and cover-triviality mechanisms --- they do not
address whether every rational terrain point lifts to one of the
covers --- and what remains open is that
$S(\Q)$ has no nondegenerate points --- equivalent to Conjecture~1 by
Proposition~\ref{prop:bridge}.  Formally, the sufficiency of
Conjecture~1
(\leanline{Gap/BMThm7Transcript.lean}{1002}{thm7\_of\_conjecture1})
and the normalisation bridge are both kernel theorems; necessity of
Conjecture~1 for Theorem~7, and the nonexistence of routes outside this
survey, are not claimed.  We make no assertion that the union of the
cover images fails to contain $S(\Q)$; no conclusion about
Conjecture~1 follows from the covers alone.  Every route to Theorem~7
in the survey above passes through that open statement.

%% ================================================================
%% §7 FORMALIZATION (5–7p)
%% SOURCE: RNT §6 + 新規 + MEASUREMENTS.md (定量表)
%% ================================================================
\section{The formalization}\label{sec:formal}

\subsection{Architecture and division of labour}\label{sec:formal:arch}

The development was produced under a fixed protocol separating three
roles.  \emph{Statements and strategy} were designed, checked, and
frozen by the author before any proof search: every definition (the
gates as derivative discriminants, the curve as a Weierstrass equation,
the $a$-map, the chart quantities, the anchor values) was cross-checked
in Sage and Magma against independent reimplementations, and the frozen
statement files were hashed.  \emph{Tactic-level proof filling} was
then delegated to large-language-model assistants (Anthropic Claude,
including its coding agent), which also drafted computer-algebra
scripts and portions of this manuscript.  \emph{Verification} was never
delegated: every Lean proof was re-verified by an independent run of
the Lean kernel, including its per-declaration axiom report, in a
process separate from the assistant runs that produced the proofs, and
again by continuous integration on independent hardware; all load-bearing
computer-algebra identities and numerical curve data were checked in
both Magma and Sage (with PARI and the LMFDB as additional cross-checks
where applicable); and no AI-generated mathematical claim was accepted
without independent human checking and, where applicable, kernel or
computer-algebra verification.  The author selected, checked, and
approved every final definition and theorem statement, and takes full
responsibility for the content of this article.  Under this protocol,
every Lean-labelled claim below is certified by the pinned kernel;
every load-bearing computer-algebra claim was reproduced independently
in Magma and Sage; ordinary mathematical deductions and literature
claims ([M], [A]) were checked by the author and lie outside the
kernel/computer-algebra trust boundary --- and nothing below rests on
unverified model output.

\subsection{The artifact: pins, certificates, and continuous
verification}\label{sec:formal:artifact}

The development compiles against mathlib (toolchain
\texttt{v4.31.0-rc1}, mathlib commit \texttt{d568c8c}) and is archived
at the frozen release \repotag{} (commit \texttt{\repocommit}) of
\url{\repourl}, with concept DOI \texttt{\conceptdoi}.  Every
declaration cited in this paper is hyperlinked to its file and line
\emph{at that tag}, so the links are permanent; the ten pins listed in
the repository README are re-resolved mechanically against the actual
tag content on every push (a fail-closed check: a pin that no longer
matches fails the build).  The dependency graph, with one node per formal declaration tracked in
the blueprint and kernel status marked per node, is at
\url{\blueprinturl}.

All computer-algebra-derived certificates --- the
\texttt{linear\_combination} cofactors of
\S\S\ref{sec:repair:backward}--\ref{sec:repair:forward}, the
elimination certificates, and the Nullstellensatz cofactors of the sign
cover --- are reproduced \emph{verbatim} in the Lean source, so their
correctness is checked by the kernel rather than trusted; the
generating scripts remain in \texttt{scripts/} as the discovery record.
Continuous integration runs three independent jobs on every push: the
Lean build with a named-declaration check (all 137 blueprint
declarations must exist) and an axiom assertion for the main theorems;
a computer-algebra re-execution job that reruns the terrain
cross-check script in a digest-pinned container and diffs its 42
checks against the committed log; and a tag self-verification job that rebuilds the
release tag itself.

\subsection{The axiom footprint, and the history it records}
\label{sec:formal:axioms}

The kernel reports the axiom footprint of
\leanline{Thm7Prime/Master.lean}{619}{thm7prime}, of the transcript
theorems of Part~I, and of the boundary theorems of Part~III as exactly
\texttt{[propext, Classical.choice, Quot.sound]} --- the standard
axioms of classical mathlib reasoning, with no custom axiom and no
\texttt{sorry}.  The development retains, deliberately, one disclosed
axiom (\leanline{Thm7Prime/Master.lean}{320}{rank\_E576i2}: the
Mordell--Weil rank of \texttt{576i2} is $1$, verified unconditionally
in two independent systems) together with the auxiliary declarations
that consume it.  These are the trace of an earlier \emph{conditional}
architecture: the forward direction of Theorem~$7'$ was first proved
under the rank hypothesis, and the unconditional reconstruction of
\S\ref{sec:repair:forward} superseded it.  The axiom was not deleted,
because the kernel's per-declaration axiom report then documents the
history of the proof as well as its final status: a reader can verify,
from the report alone, both that the main theorem is unconditional and
that it was not always so.  Mathlib's lack of a Mordell--Weil rank
theory (\S\ref{sec:open}) is thus visible in the artifact twice ---
once as the wall that stands for Conjecture~1, and once as the pressure
that produced the rank-free proof; the disclosed axiom marks the exact
point where the pressure was released.  An archived copy of the
pre-merge master file, carrying the same axiom, is retained under
\texttt{archive/} as provenance and is excluded from the live build.

\subsection{Quantitative summary}\label{sec:formal:quant}

\begin{center}
\small
\begin{tabular}{@{}lrrrrr@{}}
\toprule
component & files & lines & theorems & definitions & axioms \\
\midrule
\texttt{Gap/} (audit + boundary) & 9 & 3{,}271 & 148 & 53 & 0 \\
\texttt{Thm7Prime/} (repair) & 3 & 1{,}076 & 61 & 27 & 1 \\
\midrule
live development & 12 & 4{,}347 & 209 & 80 & 1 \\
\texttt{archive/} (provenance) & 3 & 783 & 40 & 30 & (1) \\
\bottomrule
\end{tabular}
\end{center}

\noindent Declaration counts are top-level \texttt{theorem} and
\texttt{def}-family declarations at the frozen tag; the parenthesised
axiom is the archived copy of the same \texttt{rank\_E576i2}.  The
named-declaration check covers 137 blueprint declarations.  The
blueprint itself runs to 1{,}378 lines of \LaTeX{} against 4{,}347
lines of live Lean, a ratio of roughly $1:3.2$.  The verification
scripts comprise 29 files across four suites (Magma, Sage, rank,
terrain), with a certificate index in
\texttt{scripts/CERTIFICATES.md}.
With the pinned mathlib pre-built, a full kernel verification of the
development (all twelve files, \texttt{lake build}) completes in 35
seconds of wall time (88 seconds of CPU time) on an Apple M4 Pro with
48\,GB of memory; re-checking the statement file
\texttt{Thm7Prime/Master.lean} alone takes 10 seconds.  This manuscript
runs to 2{,}052 lines of \LaTeX{} against the same 4{,}347 lines of Lean,
a ratio of roughly $1:2.1$.

\subsection{Design decisions}\label{sec:formal:design}

Three decisions shaped the development and are recorded for reuse.
First, the \emph{gate form}: rational-derivedness was formalized
through division-free discriminant gates rather than through
\texttt{Polynomial.roots} directly, because polynomial identities under
\texttt{linear\_combination} are where the kernel is strong; the cost
--- a possible statement-level slip between the gate form and the
natural form --- was then paid down by proving the equivalence in the
kernel (Theorem~\ref{thm:fidelity}).  Second, \emph{value
specialisation over sign arguments}: since $R_q$ and $S_q$ are positive
semi-definite on the chart (\S\ref{sec:boundary:integral}), negativity
certificates were unavailable for them in principle, and non-square
\emph{values} at rational points carried the exclusions instead; the
sign-based route and the value-based route coexist in the eight
verdicts, each where it applies.  Third, \emph{abstract lemmas over
base change}: the layer lemmas of \S\ref{sec:boundary:covers} were
stated over an arbitrary characteristic-zero field from the outset, so
that the geometric ($\overline{\Q}$) applications are fresh
instantiations of kernel theorems rather than informal base-change
arguments --- the kernel boundary stays exactly where the paper says it
is.

%% ================================================================
%% §8 OPEN SECTOR AND FUTURE WORK (~1.5p)
%% ================================================================
\section{The open sector}\label{sec:open}

Following the practice of stating unproved propositions as formal
definitions --- as \cite{LoefflerStoll} do for the Riemann hypothesis
--- the outstanding statement is displayed here in the exact form it
takes in the development
(\leanline{Gap/BMThm7Transcript.lean}{996}{Conjecture1\_normal};
reproduced verbatim, deliberately without proof and without
\texttt{sorry} --- introduced as a definition of a proposition rather
than as a theorem: the development contains no proof term inhabiting
it, so the source itself keeps the syntactic distinction between
stating the conjecture and proving it):

\begin{lstlisting}
def Conjecture1_normal (K : Type*) [Field K] : Prop :=
  ∀ a b : K, DistinctRoots a b → ¬ KDerived (bmQuartic a b)
\end{lstlisting}

\noindent One kernel theorem hangs from it
(\leanline{Gap/BMThm7Transcript.lean}{1002}{thm7\_of\_conjecture1}):
if Conjecture~1 (normalised) holds, then the promised conclusion T0 of
\S2.2.4 holds --- vacuously, since the conjecture empties the
hypothesis class.  This one line is the only route to T0 represented in the formal
dependency graph;
what it proves is sufficiency, not necessity.  By the normalisation
bridge (Proposition~\ref{prop:bridge}), Conjecture~1 is equivalent to
the nonexistence of nondegenerate rational points on the splitting
surface $S$; Part~III records how far computation shuts the known
approaches, and what a proof would have to decide.

One standing obstacle is a property of the current library, not of the
mathematics: at the pinned commit, mathlib does not provide the machinery needed
here for Mordell--Weil rank, Selmer computations, Chabauty's method,
or Coleman integration, so any completion of the conjecture that proceeds
through the arithmetic of elliptic surfaces cannot at present be
carried to a kernel-checked proof without either substantial new
formalization or disclosed axioms encapsulating the computer-algebra
rank computations.  For the repeated-root sector this wall was
\emph{bypassed} rather than breached --- Theorem~$7'$ is rank-free ---
and it remains standing exactly for Conjecture~1.  The reopening
condition is explicit: the required machinery must be formalised,
either in mathlib or locally.
Dated computational records (measured fibres, descent runs in
progress, the survey of cover mechanisms) are kept as research notes at
the repository root, separated from the durable blueprint.

Conjecture~1 itself remains open --- this development maps its
boundary, it does not resolve it.

%% ================================================================
%% APPENDIX
%% ================================================================
\appendix

\section{Explicit data and the formal statement}\label{app:data}

The sign-selection sextic of \S\ref{sec:repair:forward}(c) is
\begin{align*}
P(a,r,s) ={}& -15064695\,ars^4 + 9029055\,as^5 - 2983275\,rs^5\\
&+ 4985475\,s^6 + 106313580\,ars^3 - 112603320\,as^4\\
&+ 35756697\,rs^4 - 42229133\,s^5 - 143856972\,ars^2\\
&+ 415929600\,as^3 - 121880772\,rs^3 + 95245548\,s^4\\
&- 394608672\,ars - 21311424\,as^2 - 48676140\,rs^2\\
&+ 103362048\,s^3 + 833400576\,ar - 2611061568\,as\\
&+ 934428960\,rs - 758843424\,s^2 + 3732293376\,a\\
&- 1245777408\,r + 1320338880\,s - 1009822464 .
\end{align*}
The remaining explicit data (the \texttt{linear\_combination} cofactors,
the elimination certificates, and the six Nullstellensatz cofactors) are
reproduced in the Lean source and checked by the kernel; the main text
records the identity each certificate proves, and the complete
coefficient lists live in the archived development.

Following the journal's guidance that in-text code should complement
rather than reproduce the artifact, we display only the main statement;
its supporting definitions are quoted and pinned at their points of use
in \S\ref{sec:repair:statement}, and the full statement file is
\leanfile{Thm7Prime/Master.lean}.  The statements below are reproduced
verbatim from that file at the frozen tag; the correspondence between
them and Theorem~\ref{thm:main} is the statement-fidelity claim of this
paper, checkable by inspection together with
Theorem~\ref{thm:fidelity}.

\begin{lstlisting}
def Thm7Backward : Prop :=
  ∀ w z : ℚ, OnE w z → aDen w z ≠ 0 → aMap w z ≠ 0 → aMap w z ≠ 1 →
    RD211 (aMap w z)

def Thm7Forward : Prop :=
  ∀ a : ℚ, RD211 a → ∃ w z : ℚ, OnE w z ∧ aDen w z ≠ 0 ∧ aMap w z = a

def Thm7Prime : Prop := Thm7Backward ∧ Thm7Forward

theorem thm7prime : Thm7Prime := (*@$\langle$@*)thm7_backward, thm7_forward(*@$\rangle$@*)
\end{lstlisting}

%% ================================================================
%% BACK MATTER
%% ================================================================
\section*{Code availability}
{\sloppy
The Lean development, all Magma and Sage verification scripts and logs,
the blueprint, and build instructions are maintained at \url{\repourl}
and archived at the frozen release \repotag{} (commit
\texttt{\repocommit}); the release is deposited at Zenodo under concept
DOI \url{https://doi.org/\conceptdoi}.  The development descends from two earlier
archived artifacts, from which selected scripts were carried over with
a transfer manifest: the audit repository
(\url{https://doi.org/10.5281/zenodo.21465598}) and the repair
repository (\url{https://doi.org/10.5281/zenodo.21470001}).
The release revision carries the intrinsic Software Heritage identifier
\texttt{swh:1:rev:a5cf446d09b148c9c27a5232e0a0bfddd1f689ab}.
\par}

\section*{Acknowledgment}
We record with sadness that Professor R.~J. Stroeker died in October
2025, and Professor J.~A. MacDougall in May 2025, while the
verification reported here was in preparation.  The corrections
reported here are offered to the record of bodies of work that our
audit found overwhelmingly sound; we have therefore taken particular
care to distinguish typographical errors, valid intermediate results,
and the inferential gap, and to state which substantive conclusions of
the cited works remain unaffected.

\section*{Funding}
This research did not receive any specific grant from funding agencies
in the public, commercial, or not-for-profit sectors.

\section*{Declaration of competing interest}
The author declares no competing interests.

% AI 開示は §7 に一本化 (elsarticle 型の独立宣言節は AFM では不要 —
% corpus 実測では本文中の開示が標準)

\begin{thebibliography}{20}

\bibitem{BK1995}
R.~H. Buchholz, S.~M. Kelly,
On rational-derived quartics,
Bull. Austral. Math. Soc. 51 (1995) 121--132.
\url{https://doi.org/10.1017/S0004972700013940}

\bibitem{BM2000}
R.~H. Buchholz, J.~A. MacDougall,
When Newton met Diophantus: a study of rational-derived polynomials and
their extension to quadratic fields,
J. Number Theory 81 (2000) 210--233.
\url{https://doi.org/10.1006/jnth.1999.2473}

\bibitem{BC2016}
A.~Bremner, B.~Carrillo,
On $K$-derived quartics,
J. Number Theory 168 (2016) 276--291.
\url{https://doi.org/10.1016/j.jnt.2016.04.024}

\bibitem{BrascaEtAl}
A.~J. Best, C.~Birkbeck, R.~Brasca, E.~Rodriguez Boidi,
R.~Van de Velde, A.~Yang,
A complete formalization of Fermat's Last Theorem for regular primes in
Lean,
Ann. Formaliz. Math. 1 (2025) 103--132.
\url{https://doi.org/10.46298/afm.14586}

\bibitem{Carrillo2019}
B.~Carrillo,
Rational Derived Quartics and Related Topics,
Ph.D. thesis, Arizona State University, 2019.

\bibitem{Carroll1989}
C.~E. Carroll,
Polynomials all of whose derivatives have integer roots,
Amer. Math. Monthly 96 (1989) 129--130.

\bibitem{Cremona1997}
J.~E. Cremona,
Algorithms for Modular Elliptic Curves, second ed.,
Cambridge University Press, Cambridge, 1997.

\bibitem{Flynn2001}
E.~V. Flynn,
On $\Q$-derived polynomials,
Proc. Edinb. Math. Soc. (2) 44 (2001) 103--110.
\url{https://doi.org/10.1017/S0013091599000760}

\bibitem{Guy1989}
R.~K. Guy,
Unsolved problems come of age,
Amer. Math. Monthly 96, no.~10 (1989) 903--909.

\bibitem{Lean}
L.~de~Moura, S.~Ullrich,
The Lean 4 theorem prover and programming language,
in: Automated Deduction --- CADE 28, LNCS 12699, Springer, 2021,
pp. 625--635.

\bibitem{LoefflerStoll}
D.~Loeffler, M.~Stoll,
Formalizing zeta and L-functions in Lean,
Ann. Formaliz. Math. 1 (2025) 43--56.
\url{https://doi.org/10.46298/afm.15328}

\bibitem{LudwigMerten}
J.~Ludwig, C.~Merten,
Formalising the Bruhat--Tits Tree,
Ann. Formaliz. Math. 2 (2026) 55--81.
\url{https://doi.org/10.46298/afm.15738}

\bibitem{mathlib}
The mathlib Community,
The Lean mathematical library,
in: Proc. 9th ACM SIGPLAN Int. Conf. on Certified Programs and Proofs
(CPP 2020), pp. 367--381.
\url{https://doi.org/10.1145/3372885.3373824}

\bibitem{Stroeker2002}
R.~J. Stroeker,
On $\Q$-derived polynomials,
Econometric Institute Report EI 2002-30, Erasmus University Rotterdam,
2002.

\bibitem{Stroeker2006}
R.~J. Stroeker,
On $\Q$-derived polynomials,
Rocky Mountain J. Math. 36 (2006) 1705--1713.

\end{thebibliography}

\end{document}
